\documentclass[EJP]{ejpecp}

\usepackage{mathrsfs}
\usepackage{cleveref}

\emergencystretch=3em

\newtheorem{openproblem}[theorem]{Open problem}

\newcommand{\R}{\mathbb{R}}
\newcommand{\N}{\mathbb{N}}
\renewcommand{\E}{\mathbb{E}}
\newcommand{\Prob}{\mathbb{P}}
\newcommand{\Hilb}{\mathcal{H}}
\newcommand{\Dom}{\mathrm{Dom}}
\newcommand{\Ran}{\mathrm{Ran}}
\newcommand{\Energy}{\mathcal{E}}
\newcommand{\Source}{\mathsf{S}}
\renewcommand{\iota}{\mathfrak{i}}
\newcommand{\IBM}{Z}
\newcommand{\Pir}{\Pi}
\newcommand{\Channel}[1]{\mathsf{C}_{#1}}
\DeclareMathOperator{\Var}{Var}
\DeclareMathOperator{\Cov}{Cov}
\DeclareMathOperator{\supp}{supp}
\DeclareMathOperator{\sgn}{sgn}

\SHORTTITLE{Mollified It\^o--Tanaka and first-chaos dual closure for IBM}

\TITLE{A Mollified It\^o--Tanaka Formula for Iterated Brownian Motion and Its First-Chaos Dual Closure}

\AUTHORS{Ramiro~Fontes\footnote{Quijotic Research, LLC.
\EMAIL{ramirofontes@gmail.com}
\url{https://www.quijotic.com}}}

\KEYWORDS{iterated Brownian motion ; It\^o formula ; Malliavin calculus ; obstruction ; first-chaos sector ; negative Sobolev closure ; Bouleau--Yor ; bidegree atlas}

\AMSSUBJ{60H07}
\AMSSUBJSECONDARY{60J55 ; 60H05 ; 60G22}

\SUBMITTED{\today}
\ACCEPTED{}

\VOLUME{0}
\YEAR{2026}
\PAPERNUM{0}
\DOI{}

\ABSTRACT{For iterated Brownian motion $\IBM_t=X_{|Y_t|}$ with two-sided
Brownian scenery $X$ and independent Brownian clock $Y$, we prove a
mollified pathwise It\^o--Tanaka formula
\[
  f(\IBM_t^{\,\varepsilon})-f(\IBM_0^{\,\varepsilon})
   \;=\;\delta_Y\bigl(B_{f,\varepsilon}(t)\bigr)
   \;+\;C_{\IBM,\varepsilon}(f;t),
\]
holding at every $\varepsilon>0$ on the product Wiener space, with
expectation-level consistency
$\lim_{\varepsilon\to 0}\E[C_{\IBM,\varepsilon}(f;t)]=u(t,x)-f(x)$
reproducing the fourth-order Allouba--Zheng Brownian-time PDE
shadow. The unmollified pathwise correction does not survive in
classical $L^{2}(\Omega)$: under a mild nondegeneracy hypothesis,
the first-scenery-chaos projection has variance diverging as
$\varepsilon\to 0$, and the obstruction is structural. Scalar
Wick/Hermite subtraction removes the obvious diagonal
contraction but does not target the identified first-chaos
singularity; instead it introduces a different large first-chaos
contribution and disrupts the Stein cancellation. The divergent
first-chaos sector closes in
$L^{2}(\Omega_Y;H^{-1}(\R))$, with an explicit limit kernel; this
qualitative $H^{-1}$ convergence is unconditional in
$f\in\mathscr S(\R)$. As supporting structure, a fixed-$\varepsilon$
joint Wiener bidegree atlas locates the closure inside the chaos
lattice and exhibits the pure-clock $\Pir_{0,0}$ block that carries
the Brownian-time semigroup correction.}

\begin{document}

%==============================================================
\section{Introduction}\label{sec:intro}

Iterated Brownian motion
\begin{equation}\label{eq:ibm-def}
  \IBM_t\;=\;X_{|Y_t|},\qquad t\ge 0,
\end{equation}
with two-sided Brownian scenery $X=(X_s)_{s\in\R}$ and an independent
Brownian clock $Y=(Y_t)_{t\ge 0}$, is the canonical non-semimartingale
composite-noise example in continuous stochastic calculus. In its
natural filtration, $\IBM$ is not Markov and not a semimartingale;
it does not supply the finite predictable bracket structure required
by the classical semimartingale It\^o formula. Its averaged PDE
shadow, identified by Allouba--Zheng~\cite{AllouZheng}, is the
fourth-order Brownian-time equation
\begin{equation}\label{eq:az-intro}
  \partial_t u\;=\;\tfrac18\,\partial_x^{4}u
   \;+\;\frac{\partial_x^{2}f}{2\sqrt{2\pi t}}
  \qquad\text{for }u(t,x)=\E[f(x+X_{|Y_t|})],
\end{equation}
not the heat equation. The IBM observable correction operator is
therefore $\tfrac18\partial_x^{4}$ plus a boundary term, not a
second-order infinitesimal generator.

\subsection*{Contributions and structure.}
This paper proves three things.

\smallskip
\noindent\textbf{(I) A mollified pathwise It\^o--Tanaka formula
(\cref{sec:moll-ito}).}
For $\IBM_t^{\,\varepsilon}=x+X^{\varepsilon}_{|Y_t|}$ with
$X^{\varepsilon}=\rho_{\varepsilon}*X$, $f\in\mathscr S(\R)$, and
$g_{\varepsilon}(r)=f(x+X^{\varepsilon}_{r})$,
\begin{equation}\label{eq:flagship-intro}
  f(\IBM_t^{\,\varepsilon})-f(\IBM_0^{\,\varepsilon})
   \;=\;\delta_Y\bigl(B_{f,\varepsilon}(t)\bigr)
   \;+\;C_{\IBM,\varepsilon}(f;t),
\end{equation}
holds pathwise for every $\varepsilon>0$ (\cref{thm:moll-ito}), with
clock-Skorohod integrand
$B_{f,\varepsilon}(t)(s)=g_\varepsilon'(|Y_s|)\sgn(Y_s)\mathbf 1_{[0,t]}(s)$
and mollified pathwise correction
$C_{\IBM,\varepsilon}(f;t)=g_\varepsilon'(0)L_t^{0}(Y)
+\tfrac12\!\int_0^{t}g_\varepsilon''(|Y_s|)\,ds$.
In expectation, the correction reproduces the Allouba--Zheng PDE
shadow: $\lim_{\varepsilon\to 0}\E[C_{\IBM,\varepsilon}(f;t)]
=u(t,x)-f(x)$ (\cref{thm:moll-consistency}). The proofs use joint
Malliavin calculus on the product Wiener space
$\Omega_X\times\Omega_Y$.

\smallskip
\noindent\textbf{(II) The $L^{2}(\Omega)$ obstruction
(\cref{sec:obstruction}).}
Under the nondegeneracy hypothesis $(P_\eta f')(x)\not\equiv 0$ on
$(0,\infty)$, the centred first-scenery-chaos projection of
$C_{\IBM,\varepsilon}(f;t)$ has variance
\begin{equation}\label{eq:obstr-intro}
  \liminf_{\varepsilon\to 0}\,
   \Var\bigl(\Pir_1^{X}\,C_{\IBM,\varepsilon}(f;t)\bigr)
   \;=\;+\infty
\end{equation}
(\cref{thm:var-divergence}). The argument exhibits a Stein
cancellation of leading $\varepsilon^{-1}$ terms across the
two summands of $g_\varepsilon''$, leaving a residual
$\rho_\varepsilon'$-kernel whose $L^{2}$-divergence is forced by the
Bouleau--Yor spatial quadratic variation of the Brownian clock
local time. The obstruction is decisive at the $L^{2}(\Omega)$
level and is structurally Wick/Hermite-resistant: scalar
Wick-square subtraction removes the diagonal contraction in
$\Hilb_0^X\oplus\Hilb_2^X$ but does not target the $T_3$
first-chaos singularity carried by
$W_r=\partial_r^{2}X^\varepsilon\in\Hilb_1^X$, and it introduces a
large $\xi_{r,\varepsilon}$ counterterm that disrupts the Stein
cancellation (\cref{sec:wick-app}).

\smallskip
\noindent\textbf{(III) First-chaos dual closure (\cref{sec:dual}).}
The divergent first-chaos sector closes in
$L^{2}(\Omega_{Y};H^{-1}(\R))$: the conditional Malliavin kernel
$K_{\varepsilon}$ admits an interior/boundary decomposition and
converges qualitatively to
$K=L_t^{0}(Y)\,f'(x)\,\delta_0+K^{\mathrm{int}}$
(\cref{cor:combined}); this qualitative $H^{-1}$ convergence is
unconditional in $f\in\mathscr S(\R)$, while a quantitative
$\varepsilon^{1/2-\eta}$ rate is conditional on a Marcus--Rosen-type
moment assumption (\cref{assu:lt-moment,prop:interior-rate}). The
limit kernel is generically only $H^{-1}$, so the limit object is a
generalized $Y$-conditioned first-chaos functional defined by
duality against $H^{1}$ test kernels
(\cref{def:first-chaos-block}), not an ordinary $L^{2}(\Omega)$
random variable.

\smallskip
We additionally provide, for context, a fixed-$\varepsilon$
structural map of $C_{\IBM,\varepsilon}(f;t)$ across the joint
Wiener bidegree lattice $\Hilb_m^X\otimes\Hilb_n^Y$
(\cref{sec:atlas}). At fixed $\varepsilon$ the pure-clock
$\Pir_{0,0}$ sector recovers the Brownian-time semigroup correction
(\cref{thm:pure-clock}, by reference to~\cref{thm:moll-consistency});
the $(1,n)$ slices route to the dual closure block-wise
(\cref{thm:first-chaos}); and the $(2,0)$ sector admits a
six-channel energy-contraction decomposition (\cref{thm:2-0}). This
atlas is supporting structure for the main result (III), not a
fourth pillar: cross-sector $\varepsilon\to 0$ summability is open
(\cref{flag:summability}).

\subsection*{The operator-derivative bridge.}
The dual object for the singular first-chaos sector is named via a
minimal operator-derivative bridge, introduced as the opening
subsection of \cref{sec:dual} (see~\cref{sec:bridge}). The abstract
framework (operator derivative as the Hilbert adjoint of the
divergence on an energy Hilbert space, with the chain rule
following from a cylindrical reduction) is developed
in~\cite{Fontes-OperatorDerivative} and used only contextually
here. The first-chaos dual closure of~\cref{sec:dual} and the
fixed-$\varepsilon$ atlas of~\cref{sec:atlas} are stated and proved
using joint Malliavin calculus on the product Wiener space; the
bridge identifies the negative-Sobolev block $H^{-1}(\R)$ as the
test/dual space, but the proofs of
\cref{cor:combined,thm:pure-clock,thm:first-chaos,thm:2-0} are
self-contained.

The paper is organised as follows. \Cref{sec:setup} fixes the joint
Wiener setup and the mollifier estimates. \Cref{sec:moll-ito}
proves the mollified It\^o--Tanaka identity and its fourth-order
expectation-level consistency (the only statement of the
expectation--semigroup link). \Cref{sec:obstruction} proves the
first-chaos $L^{2}(\Omega)$ obstruction. \Cref{sec:dual} proves the
first-chaos dual closure in $L^{2}(\Omega_Y;H^{-1}(\R))$; an opening
subsection (\cref{sec:bridge}) names the dual object via a minimal
operator-derivative bridge and records the density reframe~(F3).
\Cref{sec:atlas} maps the
bidegree atlas at fixed $\varepsilon>0$. \Cref{sec:open} collects
the three open problems and directions for future work; two
appendices record the Wick/Hermite analysis and the Fourier-mode
PDE-shadow diagnostic.

%==============================================================
\section{Setup}\label{sec:setup}

This short section fixes the joint Wiener setup, the mollifier, and
the quantitative estimates that the rest of the paper draws on.

\subsection*{Source data, mollifier, and joint chaos.}
Let $X=(X_z)_{z\in\R}$ be a two-sided Brownian scenery with
variance $\Var(X_z)=|z|$ and covariance
$\Cov(X_z,X_{z'})=\tfrac12(|z|+|z'|-|z-z'|)$, and let
$Y=(Y_t)_{t\ge 0}$ be an independent one-sided Brownian clock with
$Y_0=0$. We write $\mathcal F_X=\sigma(X)$ and
$\mathcal F_Y=\sigma(Y)$ for the marginal $\sigma$-algebras,
$\Omega_X$ and $\Omega_Y$ for the corresponding marginal probability
spaces, and $\E_X$, $\E_Y$ for the partial expectations integrating
out the respective Gaussian factor. We fix a test function
$f\in\mathscr S(\R)$ and a basepoint $x\in\R$ throughout.

Fix a mollifier $\rho\in C_c^\infty(\R)$ that is even and
non-negative with $\int\rho=1$ and $\supp\rho\subseteq[-c,c]$, and
set $\rho_\varepsilon(\cdot)=\varepsilon^{-1}\rho(\cdot/\varepsilon)$.
The \emph{mollified scenery} is $X^\varepsilon\coloneqq\rho_\varepsilon*X$
and the \emph{mollified IBM} is
\begin{equation}\label{eq:moll-scenery}
  \IBM_t^{\,\varepsilon}\;\coloneqq\;x+X^\varepsilon_{|Y_t|}.
\end{equation}
We write $U_r\coloneqq X^\varepsilon_r$ for the value of the
mollified scenery at the time-shift $r$, and
$V_r\coloneqq\partial_r U_r=(\rho_\varepsilon'*X)(r)$ and
$W_r\coloneqq\partial_r^2 U_r=(\rho_\varepsilon''*X)(r)$ for its
first and second spatial derivatives. The pointwise variance of
$U_r$ is denoted $\sigma_\varepsilon^2(r)\coloneqq\Var(U_r)$, and
the variance of $V_r$ is $\sigma_V^2\coloneqq\Var(V_r)$, independent
of $r$ in the bulk by~\cref{lem:moll-eps-inv}. For the test
function $f$ at basepoint $x$ we set
$g_\varepsilon(r)\coloneqq f(x+U_r)$, with derivatives in $r$ given
by the chain rule:
\begin{equation}\label{eq:gprime}
  g_\varepsilon'(r)=f'(x+U_r)V_r,
  \qquad
  g_\varepsilon''(r)=f''(x+U_r)V_r^{2}+f'(x+U_r)W_r.
\end{equation}
We write $L_t^x(Y)$ for the Brownian local time of $Y$ at level $x$
up to time $t$, and $L_t^x(|Y|)$ for the spatial local time of the
reflected motion $|Y|$.

The product Wiener space $\Omega_X\times\Omega_Y$ admits the joint
Wiener chaos decomposition
\begin{equation}\label{eq:joint-chaos}
  L^{2}(\Omega_X\times\Omega_Y)
  \;=\;\bigoplus_{m,n\ge 0}\Hilb_m^X\otimes\Hilb_n^Y,
\end{equation}
where $\Hilb_m^X\subset L^2(\Omega_X)$ and
$\Hilb_n^Y\subset L^2(\Omega_Y)$ are the $m$-th scenery and $n$-th
clock Wiener chaoses. We write $\Pir_m^X$, $\Pir_n^Y$ for the
orthogonal projections onto $\Hilb_m^X$ and $\Hilb_n^Y$, extended
to the joint space as $\Pir_m^X\otimes\mathrm{Id}_Y$ and
$\mathrm{Id}_X\otimes\Pir_n^Y$ respectively. These commute, and we
write $\Pir_{m,n}=\Pir_m^X\Pir_n^Y$ for the bidegree-$(m,n)$
projection. Two identities are used below: $\Pir_0^Y=\E_Y$ is the
conditional expectation that integrates out the clock, and
$\Pir_0^X=\E_X$ integrates out the scenery.

The partial Malliavin derivatives $D_X$, $D_Y$ act on each Gaussian
factor; their adjoints are the Skorohod divergences $\delta_X=D_X^*$
and $\delta_Y=D_Y^*$~\cite{Nualart,Janson}. The scenery $X$ admits
the white-noise representation $X_z=\int_\R\sigma_z(w)\,d\xi(w)$
with the kernel
$\sigma_z(w)=\sgn(z)\mathbf 1_{[0,z]_*}(w)$. Differentiating under
the convolution gives the Malliavin derivative kernels of the
mollified scenery,
\begin{equation}\label{eq:Dw}
  D_w U_r=\xi_{r,\varepsilon}(w),
  \quad
  D_w V_r=\rho_{\varepsilon}(r-w),
  \quad
  D_w W_r=\rho_{\varepsilon}'(r-w),
\end{equation}
where $\xi_{r,\varepsilon}(w)\coloneqq\int\rho_\varepsilon(r-z)\sigma_z(w)\,dz$.

\subsection*{Quantitative estimates.}
We collect the explicit second-moment bounds used throughout, and
record the bulk covariance identities that drive the Stein
cancellation underlying both the obstruction (\cref{sec:obstruction})
and the dual closure (\cref{sec:dual}). All constants depend on
$\rho$ only via $c$, $\|\rho\|_\infty$, $\int|v|\rho(v)\,dv$,
$\int\rho^{2}$, $\int(\rho')^{2}$.

\begin{lemma}[Scenery mollifier $L^{2}$-rate, uniform in $x$]
\label{lem:moll-rate}
For every $x\in\R$,
\begin{equation}\label{eq:moll-rate-L2}
  \E\bigl[(X^{\varepsilon}_x-X_x)^{2}\bigr]\;\le\;C_\rho\,\varepsilon,
  \qquad C_\rho\;\coloneqq\;\int_\R|v|\rho(v)\,dv,
\end{equation}
hence $\E|X^\varepsilon_x-X_x|\le C_\rho^{1/2}\varepsilon^{1/2}$
uniformly in $x$. In particular $X^\varepsilon_0\to 0$ in $L^{2}$ as
$\varepsilon\to 0$ since $X_0=0$ a.s.
\end{lemma}

\begin{proof}
For two-sided Brownian scenery, $\E[(X_z-X_x)^{2}]=|z-x|$ for every
$z,x\in\R$: when $z,x$ have the same sign or one is zero, this is
the standard one-sided Brownian increment variance; when $z>0>x$,
independence of $X_z$ and $X_x$ gives
$\E[(X_z-X_x)^{2}]=\E X_z^{2}+\E X_x^{2}=z+|x|=z-x=|z-x|$ (and
symmetrically for $x>0>z$). Writing
$X^\varepsilon_x-X_x=\int_\R\rho_\varepsilon(x-z)(X_z-X_x)\,dz$ and
applying Jensen with the probability measure
$\rho_\varepsilon(x-z)\,dz$,
\(
  \E[(X^\varepsilon_x-X_x)^{2}]
  \le\int_\R\rho_\varepsilon(x-z)|z-x|\,dz
  =\int_\R\rho(v)|v|\,\varepsilon\,dv=C_\rho\varepsilon,
\)
uniform in $x$.
\end{proof}

\begin{lemma}[Scenery-derivative kernel: $L^{2}$ bound and rate]
\label{lem:moll-H}
With $H_\varepsilon(y,\xi)\coloneqq(D_X X^\varepsilon_y)(\xi)
=(\rho_\varepsilon*\sigma_\cdot(\xi))(y)$:
\begin{enumerate}
\item[(a)] $\|H_\varepsilon(y,\cdot)\|_{L^{2}(\R)}^{2}\le|y|+C_\rho\varepsilon$
  for every $y\in\R$.
\item[(b)] $\|H_\varepsilon(y,\cdot)-\sigma_y\|_{L^{2}(\R)}^{2}\le C'_\rho\varepsilon$
  for every $y\in\R$, with $C'_\rho$ depending only on $\supp\rho$.
\end{enumerate}
\end{lemma}

\begin{proof}
(a) Jensen and $\|\sigma_z\|_{L^{2}(\R)}^{2}=|z|$ give
$\int_\R H_\varepsilon(y,\xi)^{2}\,d\xi\le
\int_\R\rho_\varepsilon(y-z)|z|\,dz\le|y|+C_\rho\varepsilon$.
(b) The difference vanishes off the $O(\varepsilon)$-neighbourhood of
the jump locus of $\sigma_y$ and is uniformly bounded by $1$ there;
the $L^{2}(\R,d\xi)$ mass is $O(\varepsilon)$ uniformly in $y$.
\end{proof}

\begin{lemma}[Spatial derivative of the mollified scenery]
\label{lem:moll-eps-inv}
For $|y|>c\varepsilon$,
$V_y=(\rho_\varepsilon'*X)(y)=\int_\R\rho_\varepsilon(y-z)\,dX_z$
is centred Gaussian with
\begin{equation}\label{eq:moll-eps-inv}
  \Var(V_y)\;=\;\frac{c_\rho^{(1)}}{\varepsilon},
  \qquad c_\rho^{(1)}\coloneqq\int_\R\rho(v)^{2}\,dv.
\end{equation}
Jensen and~\eqref{eq:moll-eps-inv} extend to give
$\E[V_y^{2}]\le c_\rho^{(1)}/\varepsilon$ uniformly in $y\in\R$.
Similarly, in the bulk
$\Var(W_y)=\int_\R\rho_\varepsilon'(y-z)^{2}\,dz=c_\rho^{(2)}\varepsilon^{-3}$
with $c_\rho^{(2)}=\int(\rho')^{2}$.
\end{lemma}

\begin{proof}
For $|y|>c\varepsilon$, $z\mapsto\rho_\varepsilon(y-z)$ is supported
entirely on one side of $0$, so $\partial_y X^\varepsilon(y)=
\int\rho_\varepsilon(y-z)\,dX_z$ is an It\^o integral of an
$L^{2}$ integrand; the It\^o isometry yields
$\Var(V_y)=\|\rho_\varepsilon\|_{L^{2}}^{2}=\int\rho^{2}/\varepsilon$.
The boundary layer $|y|\le c\varepsilon$ is handled by Jensen on
the same kernel.
\end{proof}

Finally, for $|r|>c\varepsilon$ the bulk covariance identities
\begin{equation}\label{eq:bulk-cov}
  (\sigma_\varepsilon^{2})'(r)=1,\qquad
  \Cov(U_r,V_r)=\tfrac12,\qquad
  \sigma_V^{2}+\Cov(U_r,W_r)=0
\end{equation}
hold, the last by differentiating $2\Cov(U_r,V_r)=
(\sigma_\varepsilon^{2})'(r)=1$ once more; we use these without
re-proof. The boundary-layer corrections in $r\in[0,c\varepsilon]$
are quantified in~\cref{sec:dual} (\cref{lem:bdy-cov-zero,lem:bdy-layer-vanish}).

%==============================================================
\section{The mollified pathwise It\^o--Tanaka formula}\label{sec:moll-ito}

With the joint Wiener setup of~\cref{sec:setup} in hand, we now
prove the pathwise It\^o--Tanaka identity for the mollified IBM and
identify its expectation-level fourth-order consistency. We begin
with the deterministic Brownian-time semigroup correction, which
the mollified correction must reproduce in expectation as
$\varepsilon\to 0$; the mollified pathwise identity follows by
applying the classical It\^o--Tanaka formula to
$g_\varepsilon(|Y_t|)$ in the Brownian-time variable; and the
expectation-level consistency assembles the two. The
expectation-level statement (\cref{thm:moll-consistency}) is the
only locus in the paper that links the pathwise correction to the
Brownian-time semigroup; the atlas of \cref{sec:atlas} refers back
to it.

\subsection*{The Brownian-time semigroup correction.}

Before the pathwise formula, we record the expectation-level
correction identified by Allouba--Zheng. The mollified pathwise
identity will reproduce this correction in the $\varepsilon\to 0$
limit (\cref{thm:moll-consistency} below).

\begin{theorem}[Brownian-time semigroup correction]
\label{thm:bt-correction}
For $f\in\mathscr S(\R)$, set $u(t,x)\coloneqq\E[f(x+X_{|Y_t|})]$.
Then $u\in C^{\infty}((0,\infty)\times\R)$, $u(0+,x)=f(x)$, and for
every $t>0$,
\begin{equation}\label{eq:bt-correction}
  \partial_t u(t,x)
  \;=\;\frac{1}{8}\,\partial_x^{4}u(t,x)
       \;+\;\frac{1}{2\sqrt{2\pi t}}\,\partial_x^{2}f(x).
\end{equation}
\end{theorem}

\begin{proof}
Conditioning on $|Y_t|$ and using independence of $X$ and $Y$,
$u(t,x)=\int_0^\infty(P_s f)(x)\,p_t(s)\,ds$ with $P_s$ the heat
semigroup of $X$ (variance $s$, so
$\partial_s P_s f=\tfrac12\partial_x^{2}P_s f$) and
$p_t(s)=\sqrt{2/(\pi t)}\,e^{-s^{2}/(2t)}$ the half-normal density,
satisfying $\partial_t p_t=\tfrac12\partial_s^{2}p_t$ pointwise for
$s>0$, $\partial_s p_t(0)=0$, $p_t(0)=\sqrt{2/(\pi t)}$. Hence
\[
  \partial_t u(t,x)
  \;=\;\tfrac12\!\int_0^\infty\!(P_s f)(x)\,\partial_s^{2}p_t(s)\,ds.
\]
Integrate by parts twice in $s$; boundary terms at $s=\infty$ vanish:
\[
\begin{aligned}
  \tfrac12\!\int_0^\infty\!(P_s f)(x)\,\partial_s^{2}p_t(s)\,ds
  &= -\tfrac12\!\int_0^\infty\!\partial_s(P_s f)(x)\,\partial_s p_t(s)\,ds
   \quad(\partial_s p_t(0)=0)\\
  &= -\tfrac14\!\int_0^\infty\!\partial_x^{2}P_s f(x)\,\partial_s p_t(s)\,ds\\
  &= -\tfrac14\bigl[\partial_x^{2}P_s f(x)\,p_t(s)\bigr]_0^\infty
   +\tfrac14\!\int_0^\infty\!\partial_s\partial_x^{2}P_s f(x)\,p_t(s)\,ds\\
  &= \tfrac14\,\partial_x^{2}f(x)\,p_t(0)
   +\tfrac18\!\int_0^\infty\!\partial_x^{4}P_s f(x)\,p_t(s)\,ds.
\end{aligned}
\]
The boundary term equals
$\partial_x^{2}f(x)/(2\sqrt{2\pi t})$; the remaining integral equals
$\tfrac18\partial_x^{4}u(t,x)$. Smoothness on
$(0,\infty)\times\R$ and the initial condition follow from
standard regularising properties of $G_s$ and $p_t$.
\end{proof}

\begin{definition}[Semigroup correction density and integrated correction]
\label{def:sg-correction}
For $f\in\mathscr S(\R)$ and $t>0$,
\begin{equation}\label{eq:sg-density}
  \mathcal C_\IBM^{\mathrm{sg}}(f;t,x)
  \;\coloneqq\;\frac{1}{8}\,\partial_x^{4}u(t,x)
   \;+\;\frac{1}{2\sqrt{2\pi t}}\,\partial_x^{2}f(x),
  \quad
  C_\IBM^{\mathrm{sg}}(f;t,x)
  \;\coloneqq\;\int_0^t\mathcal C_\IBM^{\mathrm{sg}}(f;s,x)\,ds.
\end{equation}
The integrand is integrable on $(0,t]$:
$\int_0^t s^{-1/2}\,ds=2\sqrt t<\infty$.
\end{definition}

\begin{corollary}[Integrated semigroup form]\label{cor:sg-integrated}
For $f\in\mathscr S(\R)$ and $t>0$,
$u(t,x)-f(x)=C_\IBM^{\mathrm{sg}}(f;t,x)$.
\end{corollary}

\begin{proof}
Integrate \eqref{eq:bt-correction} from $0$ to $t$ and use
$u(0+,x)=f(x)$.
\end{proof}

\subsection*{The mollified pathwise It\^o--Tanaka identity.}

We turn now to the pathwise formula at every $\varepsilon>0$.
Conditionally on the scenery $X$, the function $r\mapsto
g_\varepsilon(r)=f(x+X_r^\varepsilon)$ is smooth, so the classical
It\^o--Tanaka formula applied to $g_\varepsilon(|Y_t|)$ in the
Brownian-time variable yields the identity below.

\begin{proposition}[Mollified pathwise It\^o--Tanaka identity]
\label{prop:moll-tanaka}
For every $\varepsilon>0$, $t\ge 0$, $f\in\mathscr S(\R)$,
\begin{equation}\label{eq:moll-tanaka}
  f(\IBM_t^{\varepsilon})-f(\IBM_0^{\varepsilon})
  \;=\;M_{\varepsilon,t}^{Y}\;+\;g_\varepsilon'(0)\,L_t^{0}(Y)
       \;+\;\tfrac12\!\int_0^t g_\varepsilon''(|Y_s|)\,ds,
\end{equation}
where
$M_{\varepsilon,t}^{Y}=\int_0^t g_\varepsilon'(|Y_s|)\sgn(Y_s)\,dY_s$
is a continuous $L^{2}$-martingale in $Y$ (conditional on $X$).
\end{proposition}

\begin{proof}
Tanaka's formula gives $d|Y_t|=\sgn(Y_t)\,dY_t+dL_t^{0}(Y)$ and
$\langle|Y|\rangle_t=t$. For each $\omega_X$, $g_\varepsilon\in
C^\infty(\R)$, so the It\^o--Tanaka formula applied to
$g_\varepsilon(|Y_t|)$ yields, using
$g_\varepsilon'(|Y_s|)\,dL_s^{0}=g_\varepsilon'(0)\,dL_s^{0}$
(since $dL_s^{0}$ is supported on $\{Y_s=0\}$),
\[
  g_\varepsilon(|Y_t|)-g_\varepsilon(0)
  =\!\int_0^t\!\! g_\varepsilon'(|Y_s|)\sgn(Y_s)\,dY_s
   +g_\varepsilon'(0)L_t^{0}(Y)
   +\tfrac12\!\int_0^t\!\! g_\varepsilon''(|Y_s|)\,ds.
\]
The left-hand side equals $f(\IBM_t^{\varepsilon})-f(\IBM_0^{\varepsilon})$
since $\IBM_t^{\varepsilon}=x+X^\varepsilon_{|Y_t|}$ and $|Y_0|=0$.
The explicit formulas for $g_\varepsilon'$, $g_\varepsilon''$
in~\eqref{eq:gprime} are the chain rule.
\end{proof}

Define the \emph{mollified pathwise correction}
\begin{equation}\label{eq:moll-correction}
  C_{\IBM,\varepsilon}(f;t)
  \;\coloneqq\;
  g_\varepsilon'(0)\,L_t^{0}(Y)\;+\;\tfrac12\!\int_0^t\!\!
   g_\varepsilon''(|Y_s|)\,ds.
\end{equation}

\begin{theorem}[Mollified IBM It\^o formula]
\label{thm:moll-ito}
For every $\varepsilon>0$, $t\ge 0$, and $f\in\mathscr S(\R)$:
\begin{enumerate}
\item (Pathwise identity)
\[
  f(\IBM_t^{\varepsilon})-f(\IBM_0^{\varepsilon})
  \;=\;\delta_Y\bigl(B_{f,\varepsilon}(t)\bigr)
       \;+\;C_{\IBM,\varepsilon}(f;t),
\]
with $B_{f,\varepsilon}(t)(s)=g_\varepsilon'(|Y_s|)\sgn(Y_s)\mathbf 1_{[0,t]}(s)$
and $C_{\IBM,\varepsilon}$ as in~\eqref{eq:moll-correction}.
\item (Skorohod--It\^o identification) The integrand
$B_{f,\varepsilon}(t)\in L^{2}(\Omega;L^{2}(\R_{+}))$, the clock
Skorohod divergence on the $Y$-adapted integrand coincides with the
It\^o integral $M_{\varepsilon,t}^{Y}$, and
$\E[\delta_Y(B_{f,\varepsilon}(t))]=0$.
\end{enumerate}
\end{theorem}

\begin{proof}
Part (1) is \cref{prop:moll-tanaka} repackaged. For (2),
$B_{f,\varepsilon}(t)$ is $Y$-adapted with continuous paths in $s$.
By \cref{lem:moll-eps-inv},
$\E[V_r^{2}]\le c_\rho^{(1)}/\varepsilon$ uniformly in $r$, so
\[
  \E\!\int_0^{t}B_{f,\varepsilon}(t)(s)^{2}\,ds
   \;\le\;\|f'\|_\infty^{2}\,t\,c_\rho^{(1)}/\varepsilon\;<\;\infty.
\]
The source clock Skorohod divergence on $Y$-adapted $L^{2}$
integrands coincides with the It\^o integral~\cite[Prop.~1.3.4]{Nualart}.
\end{proof}

Taking expectation in~\eqref{eq:moll-tanaka} eliminates the
martingale term and the boundary local-time piece (the latter by a
reflection-symmetry argument, \cref{rem:lt-zero-mean}), leaving the
expectation-level identification of $\E[C_{\IBM,\varepsilon}]$ as
the Brownian-time semigroup increment. This is the only locus in
the paper that links the pathwise correction to the
Allouba--Zheng semigroup; the bidegree atlas of~\cref{sec:atlas}
will refer back to it, not re-prove it.

\begin{theorem}[Expectation-level consistency]
\label{thm:moll-consistency}
For every $f\in\mathscr S(\R)$, $\varepsilon>0$, and $t>0$,
\begin{equation}\label{eq:moll-exp-id}
  \E[C_{\IBM,\varepsilon}(f;t)]
   \;=\;U^\varepsilon(t,x)-v^\varepsilon(0,x),
\end{equation}
where $v^\varepsilon(r,x)=\E_X[f(x+U_r)]$ and
$U^\varepsilon(t,x)=\E[f(x+U_{|Y_t|})]
=\int_0^\infty v^\varepsilon(r,x)\,p_t(r)\,dr$. Consequently
\begin{equation}\label{eq:moll-limit}
  \lim_{\varepsilon\to 0}\E[C_{\IBM,\varepsilon}(f;t)]
   \;=\;u(t,x)-f(x)
   \;=\;C_\IBM^{\mathrm{sg}}(f;t,x),
\end{equation}
with $u(t,x)=\E[f(x+X_{|Y_t|})]$ and $C_\IBM^{\mathrm{sg}}$ from
\cref{def:sg-correction}.
\end{theorem}

\begin{proof}
Take expectation in~\eqref{eq:moll-tanaka}: $\E[M_{\varepsilon,t}^{Y}]=0$,
so $\E[f(\IBM_t^\varepsilon)-f(\IBM_0^\varepsilon)]=\E[C_{\IBM,\varepsilon}(f;t)]$.
By definition $\E[f(\IBM_t^\varepsilon)]=U^\varepsilon(t,x)$ and
$\E[f(\IBM_0^\varepsilon)]=v^\varepsilon(0,x)$, giving
\eqref{eq:moll-exp-id}. For \eqref{eq:moll-limit}: Lipschitz $f$ and
\cref{lem:moll-rate} give
$\E|f(x+X^\varepsilon_r)-f(x+X_r)|\le\|f'\|_\infty C_\rho^{1/2}\varepsilon^{1/2}$
uniformly in $r$; averaging against the law of $|Y_t|$ yields
$U^\varepsilon(t,x)\to u(t,x)$. At $r=0$, $X^\varepsilon_0\to 0$ in
$L^{2}$ (\cref{lem:moll-rate} with $X_0=0$), so
$v^\varepsilon(0,x)\to f(x)$ by Lipschitz $f$. The second equality
is \cref{cor:sg-integrated}.
\end{proof}

\begin{remark}[Local-time term is expectation-invisible]
\label{rem:lt-zero-mean}
The boundary local-time term is expectation-invisible: a
reflection-symmetry argument gives $\E[g_\varepsilon'(0)L_t^{0}(Y)]=0$
for every $\varepsilon>0$ and $t>0$. Indeed, $L_t^{0}(Y)$ is
$Y$-measurable and $g_\varepsilon'(0)$ is $X$-measurable, so by
independence $\E[g_\varepsilon'(0)L_t^{0}(Y)]=\E_X[g_\varepsilon'(0)]\,\E_Y[L_t^{0}(Y)]$;
under the reflection $X_z\leftrightarrow X_{-z}$, $X^\varepsilon_0$
is invariant (even integrand against even $\rho_\varepsilon$) while
$V_0=-\int\rho_\varepsilon'(z)X_z\,dz$ picks up a sign (odd
integrand), so $\E_X[h(X^\varepsilon_0)V_0]=0$ for every Borel $h$,
and in particular $\E_X[g_\varepsilon'(0)]=0$. The full expectation
in~\eqref{eq:moll-exp-id} is therefore carried by the interior
$g_\varepsilon''$-integral; the same vanishing identity reappears
in~\cref{thm:pure-clock} via the boundary
covariance~\cref{lem:bdy-cov-zero}.
\end{remark}

The correction is fourth-order in expectation rather than in the
clock variable. The clock It\^o formula behind~\eqref{eq:moll-tanaka}
is second order in the Brownian-time variable: it is the ordinary
It\^o--Tanaka formula applied to $g_\varepsilon(|Y_t|)$. The
coefficients $g_\varepsilon'$, $g_\varepsilon''$ are however first
and second spatial derivatives of a Brownian scenery, so the
PDE shadow of $\E[C_{\IBM,\varepsilon}]$ is fourth order in the
spatial variable: $\tfrac18\partial_x^{4}u$ is the shadow of the
$f'(x+U_r)W_r$ contribution to $g_\varepsilon''$ after averaging
over the half-normal clock distribution, and
$\partial_x^{2}f/(2\sqrt{2\pi s})$ is its half-normal boundary
correction. Classical $\langle\cdot\rangle$-tracking sees only the
second-order clock term and misses the fourth-order spatial shadow
that is the actual It\^o correction for IBM.

%==============================================================
\section{The $L^{2}(\Omega)$ obstruction}\label{sec:obstruction}

The mollified formula~\eqref{eq:flagship-intro} is an identity at
every $\varepsilon>0$, and its expectation reproduces the
Allouba--Zheng PDE shadow as $\varepsilon\to 0$. The next question
is whether the pathwise correction itself converges as
$\varepsilon\to 0$ in $L^{2}(\Omega)$. It does not: the centred
fluctuation
$\tilde C_{\IBM,\varepsilon}\coloneqq C_{\IBM,\varepsilon}-\E[C_{\IBM,\varepsilon}]$
has variance diverging as $\varepsilon\to 0$
(\cref{thm:var-divergence} below). The argument exhibits a Stein
cancellation of leading $\varepsilon^{-1}$ terms in the Malliavin
kernel and isolates a residual $\rho_\varepsilon'$-kernel whose
$L^{2}$-divergence is forced by the Bouleau--Yor spatial quadratic
variation of the Brownian clock local time. The obstruction is
decisive at the $L^{2}(\Omega)$ level and is Wick/Hermite-resistant
(\cref{sec:wick-app}).

We first record three auxiliary lemmas: the Brownian local-time
Sobolev failure that drives the divergence, a conditional Stroock
projection identity that exhibits the first-chaos Malliavin kernel,
and a Plancherel-type blow-up criterion for $L^{2}$-approximants of
non-$H^{1}$ limits. The obstruction theorem then follows.

\begin{lemma}[Brownian local-time Sobolev regularity]\label{lem:lt-sobolev}
Let $Y$ be standard Brownian motion, $t>0$, and write
$L\coloneqq L_t^{\cdot}(|Y|)$ for the spatial local time of $|Y|$
on $[0,t]$, supported in $[0,M_t]$ with
$M_t=\sup_{[0,t]}|Y_s|>0$ a.s.
\begin{enumerate}
\item For every $s<\tfrac12$, $L\in H^{s}_{\mathrm{loc}}(\R)$ a.s.
\item Fix any compact $J=[a,b]\subset(0,\infty)$. On the event
$\{M_t>b\}$ (positive $\Prob_Y$-probability), the spatial occupation
of $J$ is strictly positive, hence the Bouleau--Yor spatial quadratic
variation $\mathrm{QV}_J(L)>0$ a.s., and
$L\notin H^{1}_{\mathrm{loc}}(J)$ a.s.\ on $\{M_t>b\}$.
\end{enumerate}
\end{lemma}

\begin{proof}
By Tanaka, $|Y_t|=\int_0^{t}\sgn(Y_s)\,dY_s+L_t^{0}(Y)$ is a
continuous semimartingale.
(i) $r\mapsto L_t^{r}(|Y|)$ is $1/2^{-}$ H\"older continuous
(\cite{RevuzYor}, Ch.~VI, Theorem VI.1.7 for $Y$, applied via
Tanaka to $|Y|$), so $L\in H^{s}_{\mathrm{loc}}(\R)$ for every
$s<1/2$.
(ii) The occupation identity
$\int_J L_t^{r}(|Y|)\,dr=\int_0^{t}\mathbf 1_J(|Y_s|)\,ds>0$ on
$\{M_t>b\}$ is the standard occupation-time formula
(\cite{RevuzYor}, Ch.~VI). Bouleau--Yor~\cite{BouleauYor} applied to
the continuous semimartingale $|Y|$ yields $\mathrm{QV}_J(L)>0$ a.s.\
on the same event. If $L\in H^{1}_{\mathrm{loc}}(J)$, $L$ is
absolutely continuous on $J$ with $L'\in L^{2}(J)$, and
Cauchy--Schwarz gives the spatial quadratic variation along
refining partitions bounded by $|\Delta|\,\|L'\|_{L^{2}(J)}^{2}\to 0$;
combined with $\mathrm{QV}_J(L)>0$, this is a contradiction.
\end{proof}

\begin{lemma}[Conditional first-scenery-chaos projection]
\label{lem:proj1-X}
For $F\in\mathbb D^{1,2}(\Omega_X\times\Omega_Y)$ in the scenery
direction,
\begin{equation}\label{eq:proj1-X}
  \Pir_1^{X}F
   \;=\;\delta_X\bigl(K_F(\cdot\,;Y)\bigr),
  \qquad
  K_F(w;Y)\;\coloneqq\;\E_X[D_w^{X}F\mid Y],
\end{equation}
with $K_F(\cdot;Y)\in L^{2}(\R,dw)$ a.s.
\end{lemma}

\begin{proof}
For smooth cylindrical $F$, the section $F(\cdot,\omega_Y)$ is
smooth cylindrical on the scenery factor; Stroock on $X$
(\cite[Prop.~1.3.1]{Nualart}) gives the fibrewise identity. Taking
$\E_X[\,\cdot\mid Y]$ and applying $\delta_X$ produces
\eqref{eq:proj1-X} on the cylindrical subspace. For general
$F\in\mathbb D^{1,2}$ in the scenery direction, approximate by smooth
cylindrical $F_n$ in $\mathbb D^{1,2}$
(\cite[\S1.2]{Nualart}); $L^{2}$-contraction of conditional
expectation gives $K_{F_n}\to K_F$ in $L^{2}(\Omega;L^{2}(\R))$;
chaos orthogonality gives $\Pir_1^X F_n\to\Pir_1^X F$ in $L^{2}$;
It\^o isometry gives $\delta_X K_{F_n}\to\delta_X K_F$ in $L^{2}$;
\eqref{eq:proj1-X} passes to the limit.
\end{proof}

\Cref{lem:proj1-X} applies to $C_{\IBM,\varepsilon}(f;t)$ for each
$\varepsilon>0$: $g_\varepsilon'$, $g_\varepsilon''$ are smooth
functionals of the finite Gaussian collection $(U_r,V_r,W_r)$, and
the Schwartz bounds on $f$ together with Fubini over $s\in[0,t]$
place $C_{\IBM,\varepsilon}\in\mathbb D^{1,2}$ in the scenery
direction.

\begin{lemma}[Sobolev gradient blow-up under $L^{2}$ approximation]
\label{lem:H1-blowup}
Let $u_\varepsilon\to u$ in $L^{2}(\R)$, with each
$u_\varepsilon\in C^{1}(\R)\cap H^{1}(\R)$, and $u\notin H^{1}(\R)$.
Then $\liminf_{\varepsilon\to 0}\|u_\varepsilon'\|_{L^{2}(\R)}=+\infty$.
\end{lemma}

\begin{proof}
By Plancherel,
$\|u_\varepsilon'\|_{L^{2}}^{2}=(2\pi)^{-1}\int k^{2}|\widehat{u_\varepsilon}|^{2}\,dk$.
For any $R>0$, dominated convergence on $|k|\le R$ gives
$\int_{|k|\le R}k^{2}|\widehat{u_\varepsilon}|^{2}\,dk\to
\int_{|k|\le R}k^{2}|\widehat u|^{2}\,dk$. Letting $R\to\infty$ with
$\int k^{2}|\widehat u|^{2}\,dk=+\infty$ gives the conclusion.
\end{proof}

With these three inputs in place, the obstruction is now a direct
computation.

\begin{theorem}[$L^{2}$-divergence of the first-chaos projection]
\label{thm:var-divergence}
Let $f\in\mathscr S(\R)$, $x\in\R$, $t>0$. Assume the nondegeneracy
condition
\begin{equation}\label{eq:nondegen}
  C(\eta)\;\coloneqq\;\E[f'(x+X_\eta)]\;=\;(P_\eta f')(x),
  \qquad\eta>0,
\end{equation}
is not identically zero on $(0,\infty)$ (in particular, this holds
whenever $f'(x)\ne 0$, since $C(\eta)\to f'(x)$ as $\eta\to 0^{+}$).
Then
\begin{equation}\label{eq:var-diverges}
  \liminf_{\varepsilon\to 0}\,
   \Var\bigl(\Pir_1^{X}\,C_{\IBM,\varepsilon}(f;t)\bigr)\;=\;+\infty,
\end{equation}
and in particular $C_{\IBM,\varepsilon}(f;t)$ does not converge in
$L^{2}(\Omega)$ as $\varepsilon\to 0$.
\end{theorem}

\begin{proof}
Set $U_r=X^\varepsilon_r$, $V_r=\partial_rX^\varepsilon_r$,
$W_r=\partial_r^{2}X^\varepsilon_r$. In the bulk ($r>c\varepsilon$),
$\sigma_V^{2}=c_\rho^{(1)}/\varepsilon$, $\Cov(U_r,V_r)=\tfrac12$,
$\Cov(U_r,W_r)=-c_\rho^{(1)}/\varepsilon$
(via~\eqref{eq:bulk-cov}). Malliavin kernels:
$D_wV_r=\rho_\varepsilon(r-w)$, $D_wW_r=\rho_\varepsilon'(r-w)$,
$D_wU_r=\xi_{r,\varepsilon}(w)$.

\emph{Step 1: first-chaos kernel decomposition.}
By \cref{lem:proj1-X}, the conditional Malliavin kernel of
$\Pir_1^{X}C_{\IBM,\varepsilon}$ is
$K_\varepsilon(w\mid Y)=\E_X[D_w^{X}C_{\IBM,\varepsilon}\mid Y]$,
splitting along the two summands of $C_{\IBM,\varepsilon}$:
\begin{equation}\label{eq:Kdecomp}
  K_\varepsilon(w\mid Y)
   \;=\;K_\varepsilon^{\mathrm{int}}(w\mid Y)
   \;+\;K_\varepsilon^{\mathrm{bdy}}(w\mid Y),
\end{equation}
$K_\varepsilon^{\mathrm{bdy}}(w\mid Y)=L_t^{0}(Y)\,\E_X[D_wg_\varepsilon'(0)]$,
$K_\varepsilon^{\mathrm{int}}(w\mid Y)=\tfrac12\E_X[\int_0^t D_w g_\varepsilon''(|Y_s|)\,ds\mid Y]$.

\emph{Step 2: spatial support of $K_\varepsilon^{\mathrm{bdy}}$.}
From $g_\varepsilon'(0)=f'(x+U_0)V_0$,
$\E_X[D_wg_\varepsilon'(0)]=\E_X[f''(x+U_0)V_0]D_wU_0
+\E_X[f'(x+U_0)]D_wV_0$.
Both kernels are supported in $\{|w|\le c\varepsilon\}$:
$D_wV_0=\rho_\varepsilon(w)$ has compact support there; and
$D_wU_0=\xi_{0,\varepsilon}(w)=\int\rho_\varepsilon(z)\sgn(z)\mathbf 1_{w\in[0,z]_*}\,dz$
vanishes for $|w|>c\varepsilon$ because $\rho_\varepsilon(z)$ is
supported in $|z|\le c\varepsilon$. Hence
\begin{equation}\label{eq:Kbdy-support}
  \supp\,K_\varepsilon^{\mathrm{bdy}}(\cdot\mid Y)
   \;\subseteq\;[-c\varepsilon,c\varepsilon],\qquad\forall\,\varepsilon>0.
\end{equation}

\emph{Step 3: interior kernel via Stein/Malliavin.}
For $g_\varepsilon''(r)=f''(x+U_r)V_r^{2}+f'(x+U_r)W_r$,
\begin{equation}\label{eq:Dgpp}
\begin{aligned}
  \E[D_wg_\varepsilon''(r)]
   &=\bigl(\E[f'''(x+U_r)V_r^{2}]+\E[f''(x+U_r)W_r]\bigr)\xi_{r,\varepsilon}(w)\\
   &\quad+2\E[f''(x+U_r)V_r]\,\rho_\varepsilon(r-w)
   +\E[f'(x+U_r)]\,\rho_\varepsilon'(r-w).
\end{aligned}
\end{equation}
Stein for $(U,V)$ jointly Gaussian:
$\E[h(x+U)V]=\Cov(U,V)\E[h'(x+U)]$ and
$\E[h(x+U)V^{2}]=\sigma_V^{2}\E[h(x+U)]+\Cov(U,V)^{2}\E[h''(x+U)]$.
In the bulk,
$\E[f'''(x+U_r)V_r^{2}]+\E[f''(x+U_r)W_r]
=(\sigma_V^{2}+\Cov(U_r,W_r))\E[f'''(x+U_r)]+\tfrac14\E[f^{(5)}(x+U_r)]
=\tfrac14\E[f^{(5)}(x+U_r)]$, and
$2\E[f''(x+U_r)V_r]=\E[f'''(x+U_r)]$, so the two $\varepsilon^{-1}$
contributions cancel. Writing
$A_\varepsilon(r)=\tfrac14\E[f^{(5)}(x+U_r)]$,
$B_\varepsilon(r)=\E[f'''(x+U_r)]$,
$C_\varepsilon(r)=\E[f'(x+U_r)]$, with limits $A,B,C$ given by
$A(r)=\tfrac14(P_{|r|}f^{(5)})(x)$, $B(r)=(P_{|r|}f''')(x)$,
$C(r)=(P_{|r|}f')(x)$ matching~\eqref{eq:nondegen},
\begin{equation}\label{eq:Dkernel}
  \E[D_wg_\varepsilon''(r)]
  =A_\varepsilon(r)\xi_{r,\varepsilon}(w)
   +B_\varepsilon(r)\rho_\varepsilon(r-w)
   +C_\varepsilon(r)\rho_\varepsilon'(r-w).
\end{equation}
The boundary-layer corrections at $r\in[0,c\varepsilon]$ play no
role in the localised argument below. Setting
$L\coloneqq L_t^{\cdot}(|Y|)$ (supported in $[0,M_t]$), integration
against $L\,dr$ gives
\begin{equation}\label{eq:T123}
  K_\varepsilon^{\mathrm{int}}(w\mid Y)
   =\tfrac12\bigl(T_1+T_2+T_3\bigr)(w),
\end{equation}
where $T_1(w)=\int_0^{\infty}A_\varepsilon(r)L(r)\xi_{r,\varepsilon}(w)\,dr$,
$T_2(w)=(\rho_\varepsilon*B_\varepsilon L)(w)$,
$T_3(w)=-\partial_w(\rho_\varepsilon*C_\varepsilon L)(w)$.

\emph{Step 4: spatial localisation away from the boundary.}
By~\eqref{eq:nondegen} and analyticity of $C(\eta)$ on $(0,\infty)$,
the zero set of $C$ is discrete; choose $J=[a,b]\subset(0,\infty)$
with $\inf_J|C|\ge c_J>0$, and a non-empty closed subinterval
$J^{\circ}=[a',b']\subset(a,b)$. Fix smooth cutoff
$\chi\in C_c^{\infty}(\R)$ with $\chi\equiv 1$ on $J^{\circ}$,
$\supp\chi\subset J$. Restrict to $\varepsilon$ with $c\varepsilon<a/2$.
By Step~2, $\supp K_\varepsilon^{\mathrm{bdy}}\cap\supp\chi=\emptyset$,
so $\chi K_\varepsilon^{\mathrm{bdy}}\equiv 0$ and
\begin{equation}\label{eq:localised}
  \Var_X(\Pir_1^{X}C_{\IBM,\varepsilon}\mid Y)
  =\|K_\varepsilon(\cdot\mid Y)\|_{L^{2}(\R)}^{2}
  \;\ge\;\tfrac14\|\chi(T_1+T_2+T_3)\|_{L^{2}}^{2}.
\end{equation}

\emph{Step 5: $\|\chi T_1\|$ and $\|\chi T_2\|$ are bounded
uniformly.}
$|\chi|\le 1$ gives $\|\chi T_j\|\le\|T_j\|$.
$\|T_1\|^{2}\le(\|A\|_\infty t)^{2}(M_t+c\varepsilon)$ via the
covariance bound and $\int L=t$;
$\|T_2\|^{2}\le\|B\|_\infty^{2}\|L\|_\infty^{2}M_t$ by Young's
inequality. Both bounds are uniform in $\varepsilon$.

\emph{Step 6: $\|\chi T_3\|\to\infty$ on $\{M_t>b\}$.}
Set $u_\varepsilon=\rho_\varepsilon*(C_\varepsilon L)$. Heat-semigroup
time-continuity gives $C_\varepsilon\to C$ uniformly on compacts of
$(0,\infty)$; $L$ has compact support in $[0,M_t]$; so
$u_\varepsilon\to CL$ in $L^{2}(\R)$, hence $\chi u_\varepsilon\to\chi CL$
in $L^{2}(\R)$. Each $\chi u_\varepsilon$ is in $C^{1}\cap H^{1}$.

\emph{Claim: $\chi CL\notin H^{1}(\R)$ a.s.\ on $\{M_t>b\}$.}
On $J^{\circ}$, $\chi\equiv 1$ and $|C|\ge c_J>0$; if
$\chi CL\in H^{1}(\R)$, then $CL|_{J^{\circ}}\in H^{1}_{\mathrm{loc}}(J^{\circ})$,
and dividing by smooth nonvanishing $C$ gives
$L\in H^{1}_{\mathrm{loc}}(J^{\circ})$. \Cref{lem:lt-sobolev}(ii)
on $J^{\circ}\subset(0,\infty)$ yields the contradiction
$L\notin H^{1}_{\mathrm{loc}}(J^{\circ})$ a.s.\ on $\{M_t>b'\}\supseteq\{M_t>b\}$.

\Cref{lem:H1-blowup} applied to
$u_\varepsilon^{\chi}=\chi u_\varepsilon\to\chi CL\notin H^{1}(\R)$
on $\{M_t>b\}$ gives
$\liminf_{\varepsilon\to 0}\|(\chi u_\varepsilon)'\|_{L^{2}}=+\infty$
a.s.\ on $\{M_t>b\}$. Now $\chi T_3=-(\chi u_\varepsilon)'+\chi'u_\varepsilon$
with $\|\chi'u_\varepsilon\|_{L^{2}}\to\|\chi'CL\|_{L^{2}}<\infty$;
hence $\|\chi T_3\|\to+\infty$ a.s.\ on $\{M_t>b\}$.

\emph{Step 7: conclusion via Fatou.}
By the reverse triangle inequality on~\eqref{eq:localised} and
Steps~5--6,
\(
  \|K_\varepsilon(\cdot\mid Y)\|_{L^{2}(\R)}
   \ge\tfrac12(\|\chi T_3\|-\|\chi T_1\|-\|\chi T_2\|)
   \to+\infty
\)
a.s.\ on $\{M_t>b\}$, an event of positive $\Prob_Y$-probability
(since $\{Y_t>b\}\subseteq\{M_t>b\}$). \Cref{lem:proj1-X} and
Fatou under $\E_Y$ give
$\liminf_\varepsilon\E_Y\Var_X(\Pir_1^{X}C_{\IBM,\varepsilon}\mid Y)=+\infty$;
the law of total variance yields~\eqref{eq:var-diverges}.
\end{proof}

\begin{remark}[Anatomy of the obstruction]\label{rem:obstr-anatomy}
The proof exhibits two structural facts. (a) The naive
$\varepsilon^{-1}$-divergent first-chaos contributions from
$f''(x+U_r)V_r^{2}$ and $f'(x+U_r)W_r$ cancel in the bulk Malliavin
kernel~\eqref{eq:Dkernel}: this is the Stein cancellation. (b) The residual is the $T_3$ summand, a
$\rho_\varepsilon'$-kernel with coefficient $C_\varepsilon L$; its
$L^{2}(\R)$ blow-up is forced by $\chi CL\notin H^{1}(\R)$, in turn
forced by the Bouleau--Yor spatial quadratic variation of the
reflected Brownian local time. The first-chaos divergence is
therefore not an artifact of mollification but a structural feature
of the composite scenery-at-clock object.
\end{remark}

\begin{remark}[Wick/Hermite is not a stabilising renormalisation]
\label{rem:wick-not-sufficient}
The diagonal contraction
$\mathrm{Contr}_\varepsilon(r)=\E[V_r^{2}]=c_\rho^{(1)}/\varepsilon$
is the obvious renormalisation candidate, equivalent to Wick/Hermite
square ordering in the Gaussian source atlas. The subtraction
removes the even-sector diagonal contraction in $\Hilb_0^X\oplus\Hilb_2^X$,
but it neither touches the $T_3$ first-chaos singularity (carried
by $W_r=\partial_r^{2}X^\varepsilon\in\Hilb_1^X$ with
$\Var\sim\varepsilon^{-3}$) nor preserves the Stein cancellation
that isolates it: the counterterm introduced by the subtraction is
itself a large first-chaos $\xi_{r,\varepsilon}$-contribution.
Scalar Wick-square subtraction therefore does not remove the
first-chaos obstruction and cannot by itself provide the
stabilising renormalisation of the full correction. The detailed
analysis is in~\cref{sec:wick-app}.
\end{remark}

\Cref{thm:var-divergence} rules out the raw $L^{2}(\Omega)$ limit
of the unmollified pathwise correction. A natural target is the
antidual of a smooth-test subspace in a rigged Hilbert hierarchy on
the second-order Malliavin energy of the IBM source; constructing
the full dual concretely is open (\cref{flag:summability}). The
first-scenery-chaos slice of that dual is built in~\cref{sec:dual}:
the negative-Sobolev block $H^{-1}(\R)$ on the scenery variable
hosts the limit kernel of $\Pir_1^X C_{\IBM,\varepsilon}$.

%==============================================================
\section{First-chaos dual closure}\label{sec:dual}

The obstruction of~\cref{sec:obstruction} localises the divergent
content of $\Pir_1^{X}C_{\IBM,\varepsilon}$ on the scenery variable
$w\in\R$ and identifies the offending kernel as
$-\partial_w(\rho_\varepsilon*C_\varepsilon L)$, which has no
$L^{2}(\R)$ limit. Dropping one spatial derivative absorbs the
blow-up: the limit lives in $H^{-1}(\R)$. We establish this in this
section. The qualitative $H^{-1}$ convergence
(\cref{thm:interior,thm:boundary,cor:combined}) is unconditional in
$f\in\mathscr S(\R)$ for the first-chaos sector; cross-sector
$\varepsilon\to 0$ summability is open (\cref{flag:summability}).
A quantitative $\varepsilon^{1/2-\eta}$ rate
(\cref{prop:interior-rate}) requires a Marcus--Rosen-type moment
estimate~(\cref{assu:lt-moment}) on the reflected local time, stated
as a hypothesis of the proposition.

The section opens with the minimal operator-derivative bridge that
names the dual object, then proves the closure in three parts: the
kernel decomposition and its boundary-layer residual; the interior
$H^{-1}$ convergence (qualitative and conditional rate); the
boundary contribution and the combined closure.

\subsection*{The operator-derivative bridge.}\label{sec:bridge}

This subsection introduces the minimum operator-derivative structure
needed to name the dual object hosting the $\varepsilon\to 0$ limit
of $\Pir_1^X C_{\IBM,\varepsilon}$: an integration-by-parts identity
fixing notation, the negative-Sobolev block on the scenery factor
that will host the limit kernel, and the density question that the
construction reframes. The abstract operator-derivative framework
(adjointness, cylindrical reduction, chain and Leibniz rules) is
developed in~\cite{Fontes-OperatorDerivative} and cited contextually;
the convergence theorems below are stated and proved using joint
Malliavin calculus directly.

On the product Wiener space $(\Omega,\mathcal F,\Prob)$ with source
$\sigma$-algebra $\mathcal F_\Source=\mathcal F_X\vee\mathcal F_Y$
and joint Skorohod divergence
$\delta_\Source(u_X,u_Y)=\delta_X u_X+\delta_Y u_Y$ on
$\Dom(\delta_X)\oplus\Dom(\delta_Y)$, the joint Malliavin derivative
$D_\Source=(D_X,D_Y)$ is the Hilbert adjoint of $\delta_\Source$ in
the energy space $\Energy_\Source=L^{2}(\Omega;L^{2}(\R)\oplus L^{2}(\R_{+}))$,
the defining identity being the integration by parts
\begin{equation}\label{eq:bridge-ibp}
  \E[F\cdot\delta_\Source u]=\E[\langle D_\Source F,u\rangle_{\Energy_\Source}]
\end{equation}
for $F\in\mathbb D^{1,2}$ on each Gaussian factor and
$u\in\Dom(\delta_\Source)$. In the joint-Gaussian setting at hand
$D_X$, $D_Y$ are the classical Malliavin derivatives~\cite{Nualart,Janson}
and~\eqref{eq:bridge-ibp} is classical; we record it only to fix
notation.

The obstruction~\eqref{eq:var-diverges} is localised on the scenery
variable $w\in\R$ (\cref{rem:obstr-anatomy}); the natural test/dual
block is therefore the standard Sobolev pair on that factor.

\begin{definition}[First-chaos test space, dual, and generalized
first-chaos functional]\label{def:first-chaos-block}
Set
\[
 \Energy_{\IBM,1}^{+}\;\coloneqq\;L^{2}(\Omega_Y;H^{1}(\R,dw)),
 \qquad
 \Energy_{\IBM,1}^{-}\;\coloneqq\;L^{2}(\Omega_Y;H^{-1}(\R,dw)),
\]
with the fibrewise duality pairing
$\langle K(\cdot\mid Y),\Phi(\cdot\mid Y)\rangle_{H^{-1},H^{1}}$
extending the $L^{2}(\R)$ inner product on each fibre. A
\emph{generalized $Y$-conditioned first-chaos functional} is an
object specified by a kernel
$K\in\Energy_{\IBM,1}^{-}=L^{2}(\Omega_Y;H^{-1}(\R))$; it acts on
$Y$-random test integrands $\Phi\in\Energy_{\IBM,1}^{+}=L^{2}(\Omega_Y;H^{1}(\R))$
by the \emph{first-chaos pairing}
\begin{equation}\label{eq:first-chaos-pairing}
 \bigl\langle K,\Phi\bigr\rangle_{1,*}
 \;\coloneqq\;\E_{Y}\!\bigl[\langle K(\cdot\mid Y),\Phi(\cdot\mid Y)\rangle_{H^{-1},H^{1}}\bigr],
\end{equation}
finite by Cauchy--Schwarz in $L^{2}(\Omega_Y)$. A deterministic test
$\phi\in H^{1}(\R)$ is the special case $\Phi(Y,w)=\phi(w)$. When
$K(\cdot\mid Y)\in L^{2}(\R)$ a.s., the pairing reduces to the
fibrewise $L^{2}(\R)$-pairing and the functional is realised by the
ordinary first-chaos random variable
$I_1^X(K(\cdot\mid Y))\in L^{2}(\Omega)$; when $K(\cdot\mid Y)$ is
only in $H^{-1}(\R)$, no such $L^{2}(\Omega_X)$ random variable
exists, and the functional is specified by~\eqref{eq:first-chaos-pairing}
alone, on $Y$-random test kernels $\Phi$ identified fibrewise with
first-chaos test variables $I_1^X(\Phi(\cdot\mid Y))$.
\end{definition}

\begin{remark}[Why $H^{-1}$ is the natural choice]\label{rem:why-Hminus1}
Two minimal facts are available. (a) $\partial_w\colon L^{2}(\R)\to H^{-1}(\R)$
is bounded with norm $\le 1$, since for $\phi\in H^{1}(\R)$,
$|\langle\partial_w f,\phi\rangle|=|\langle f,\phi'\rangle|\le
\|f\|_{L^{2}}\|\phi\|_{H^{1}}$. (b) The offending kernel
$T_3^{\varepsilon}=-\partial_w(\rho_\varepsilon*C_\varepsilon L)$ from
the proof of \cref{thm:var-divergence} satisfies
$\rho_\varepsilon*C_\varepsilon L\to CL$ in $L^{2}(\R)$ a.s., so
$\partial_w(\rho_\varepsilon*C_\varepsilon L)\to\partial_w(CL)$ in
$H^{-1}(\R)$ a.s. Together these close the singular channel in
$H^{-1}$. The sharper local threshold is governed by Bouleau--Yor
regularity: $L\in H^{1/2-\eta}_{\mathrm{loc}}$ a.s., so
$\partial_w(CL)\in H^{-1/2-\eta}_{\mathrm{loc}}\subset H^{-1}_{\mathrm{loc}}$.
We work in $H^{-1}$ as the natural Sobolev-dual block; no claim of
optimality is made.
\end{remark}

\begin{remark}[``Quartic'' vs.\ Wiener chaos degree]
\label{rem:quartic-not-chaos}
A label sometimes attached to $\Energy_{\IBM}^{-}$ in this
literature is ``quartic dual,'' referring to the order of operator
composition that produced the kernel: the second derivative
$D_\Source^{(2)}$ acts twice in assembling $g_\varepsilon''$ from
$f$ and the scenery, and the test side requires four-fold structure.
This is not the Wiener-chaos degree of the resulting random
variable. The singular object of \cref{thm:var-divergence} is
genuinely first-$X$-chaos; the first-chaos sector is a single
contraction channel of the four-fold composite, not a four-variable
kernel space. We use this convention throughout.
\end{remark}

\smallskip
\noindent\textbf{The density question.}
\Cref{def:first-chaos-block} is on the scenery side. On the
observable side, a natural \emph{candidate} intrinsic IBM derivative
is obtained by joint-Malliavin compression on the maximal pullback
core
$\mathcal C_\IBM=\{F\in L^{2}(\mathcal F_\IBM):\iota F\in\Dom(D_X)\cap\Dom(D_Y)\}$:
write $D_\IBM F\coloneqq(D_X\iota F,D_Y\iota F)$ for $F\in\mathcal C_\IBM$,
which is closed on this core. When the core is dense in
$L^{2}(\mathcal F_\IBM)$, the candidate operator admits the usual
Hilbert adjoint $\delta_\IBM=D_\IBM^{*}$, expected to satisfy the
projected divergence formula
$\delta_\IBM u=\E[\delta_X u_X+\delta_Y u_Y\mid\mathcal F_\IBM]$ on
the lifted source-divergence domain. In the present IBM setting we
do not assert such density; accordingly, the candidate intrinsic
adjoint/divergence should be understood as part of the open
rigged-Hilbert-space problem, not as an object used below. The
abstract framework is developed
in~\cite{Fontes-OperatorDerivative}.

We would like to point out the following.

\begin{openproblem}[F3: density of $\mathcal C_\IBM$]\label{flag:density}
We do not assert density of $\mathcal C_\IBM$ in
$L^{2}(\mathcal F_\IBM)$. The construction suggests density may
fail, because the natural observables of IBM involve pointwise
evaluation of the scenery at the clock $X_{|Y_t|}$, which lies
outside the Sobolev domain of $D_X$ in any unmollified sense; the
mollified surrogates $X^{\varepsilon}_{|Y_t|}$ lie on the joint
Sobolev domain at every $\varepsilon>0$, but they are not
evidently representable as $\mathcal F_\IBM$-measurable observables,
and in any case the resulting approximating class is not shown here
to be closed under the pullback to $\mathcal C_\IBM$. Under this
scenario the natural home for the unmollified pathwise correction
$C_\IBM(f;t)=\lim_{\varepsilon\to 0}C_{\IBM,\varepsilon}(f;t)$ is
not $L^{2}(\mathcal F_\IBM)$.
\end{openproblem}

The reframe is therefore the following. The candidate intrinsic IBM
derivative $D_\IBM$ is closed on its maximal pullback core
$\mathcal C_\IBM$; we are not suggesting $\mathcal C_\IBM$ is dense
in $L^{2}(\mathcal F_\IBM)$, and we expect it is not, for the
structural reason above. The unmollified pathwise correction
$C_\IBM(f;t)$ is consequently not expected to live in
$L^{2}(\Omega)$: \cref{thm:var-divergence} rules out its first-chaos
component on classical $L^{2}$ grounds, and the natural home is a
dual / rigged-Hilbert-space object. The first-chaos slice of that
dual is what \cref{sec:dual} constructs; the full cross-sector
construction is open (\cref{flag:summability}). The label
``quartic observable energy'' used in this literature refers to the
operator-composition order of the test side, not to chaos degree,
and is not invoked structurally in any proof below.

\Cref{cor:combined} of~\cref{sec:dual} proves the corresponding
first-chaos limit in the dual topology that the obstruction selects,
which gives the reframe constructive content.
Cross-sector summability is the remaining frontier
(\cref{flag:summability}). Throughout the rest of the paper,
``operator derivative'' is invoked terminologically only, as the
bridge that names the $H^{-1}$ dual block; no chain rule,
cylindrical reduction, or second-order $D_\IBM^{(2)}$ is asserted
or needed in~\cref{sec:dual,sec:atlas}.

%==============================================================
\subsection*{The kernel decomposition.}

We re-derive the first-chaos Malliavin kernel of
$\Pir_1^{X}\,C_{\IBM,\varepsilon}(f;t)$ from first principles. Recall
$U_r=X^\varepsilon_r$, $V_r=\partial_r U_r$, $W_r=\partial_r^{2}U_r$
with white-noise representations~\eqref{eq:Dw}. We use
$c_\rho^{(1)}=\int\rho^{2}$ and $c_\rho^{(2)}=\int(\rho')^{2}$ as in
\cref{sec:setup}, together with $C_\rho=\int|v|\rho(v)\,dv$ and
$K_\rho=\tfrac12\!\iint\rho(v_1)\rho(v_2)|v_1-v_2|\,dv_1dv_2$.

The Gaussian Stein identities (re-derived):
$\E[h(U)V]=\Cov(U,V)\E[h'(U)]$;
$\E[h(U)V^{2}]=\Cov(U,V)^{2}\E[h''(U)]+\sigma_V^{2}\E[h(U)]$ for
jointly centred Gaussian $(U,V)$ and smooth $h$.

Applying $D_w$ to
$g_\varepsilon''(r)=f''(x+U_r)V_r^{2}+f'(x+U_r)W_r$ and taking
$\E_X$ gives, in the bulk $r>c\varepsilon$
where~\eqref{eq:bulk-cov} holds, the cancellation already used in
\cref{thm:var-divergence}: writing
$A_\varepsilon(r)=\tfrac14\E_X[f^{(5)}(x+U_r)]$,
$B_\varepsilon(r)=\E_X[f'''(x+U_r)]$,
$C_\varepsilon(r)=\E_X[f'(x+U_r)]$ (bounded uniformly in
$(r,\varepsilon)$ with heat-semigroup limits
$A(r)=\tfrac14(P_{|r|}f^{(5)})(x)$, $B(r)=(P_{|r|}f''')(x)$,
$C(r)=(P_{|r|}f')(x)$ as in~\eqref{eq:nondegen}),
\begin{equation}\label{eq:kernel-form}
 \E_{X}[D_{w}g_{\varepsilon}''(r)]
 =A_{\varepsilon}(r)\xi_{r,\varepsilon}(w)+B_{\varepsilon}(r)\rho_{\varepsilon}(r-w)
 +C_{\varepsilon}(r)\rho_{\varepsilon}'(r-w).
\end{equation}

\begin{lemma}[Boundary covariance vanishes]\label{lem:bdy-cov-zero}
For two-sided Brownian scenery $X$ and even mollifier $\rho$,
\(
 \Cov(U_{0},V_{0})=\tfrac12(\sigma_{\varepsilon}^{2})'(0)=0,
\)
hence $U_{0}\perp V_{0}$.
\end{lemma}

\begin{proof}
$(\sigma_{\varepsilon}^{2})'(r)=I_{1}'(r)-\tfrac12 I_{2}'(r)$ with
$I_{1}(r)=\int\rho_{\varepsilon}(r-z)|z|\,dz$ and
$I_{2}(r)=\iint\rho_{\varepsilon}(r-z_{1})\rho_{\varepsilon}(r-z_{2})|z_{1}-z_{2}|\,dz_{1}dz_{2}$.
At $r=0$, translation gives
$I_{1}'(0)=\int\rho_{\varepsilon}'(z)|z|\,dz=0$ by odd--even pairing.
By symmetry in $(z_1,z_2)$,
$I_{2}'(0)=2\!\iint\rho_{\varepsilon}'(z_{1})\rho_{\varepsilon}(z_{2})|z_{1}-z_{2}|\,dz_{1}dz_{2}$;
integration by parts in $z_1$ (compact support of $\rho_\varepsilon$
kills the boundary) gives
$I_{2}'(0)=-2\!\iint\rho_{\varepsilon}(z_{1})\rho_{\varepsilon}(z_{2})\sgn(z_{1}-z_{2})\,dz_{1}dz_{2}=0$
by antisymmetry of $\sgn(z_1-z_2)$ against the symmetric product
measure. Joint Gaussianity and vanishing covariance give
independence.
\end{proof}

\begin{remark}[Boundary kernel coefficient]
\label{rem:bdy-coeff}
The consequence for the boundary kernel is immediate. Because
$V_0\perp U_0$, $\E_X[f''(x+U_0^\varepsilon)V_0]=\E[V_0]\,\E_X[f''(x+U_0^\varepsilon)]=0$;
the $\xi_{0,\varepsilon}$ piece in $\E_X[D_w g_\varepsilon'(0)]$
therefore vanishes identically, leaving only the
$\rho_\varepsilon$-kernel with coefficient $C_\varepsilon(0)$. This
rules out the naively conjectured $\delta_0'$ term in the limit of
$K_\varepsilon^{\mathrm{bdy}}$ (see~\cref{thm:boundary}).
\end{remark}

Combining the Stein audit with the occupation-density formula
$\int_0^t h(|Y_s|)\,ds=\int_0^\infty h(r)L_t^{r}(|Y|)\,dr$
($L\coloneqq L_t^{\cdot}(|Y|)$), the conditional Malliavin kernel
of $\Pir_1^{X}\,C_{\IBM,\varepsilon}(f;t)$ is
\begin{equation}\label{eq:Keps-decomp}
 K_{\varepsilon}(w\mid Y)
 =K_{\varepsilon}^{\mathrm{bdy}}(w\mid Y)
 +\tfrac12\bigl(T_{1}^{\varepsilon}+T_{2}^{\varepsilon}+T_{3}^{\varepsilon}\bigr)(w\mid Y)
 +R_{\varepsilon}^{\mathrm{layer}}(w\mid Y),
\end{equation}
with
\begin{align}
 K_{\varepsilon}^{\mathrm{bdy}}(w\mid Y)
 &=L_{t}^{0}(Y)\cdot C_{\varepsilon}(0)\,\rho_{\varepsilon}(w),\label{eq:Kbdy}\\
 T_{1}^{\varepsilon}(w\mid Y)
 &=\int_{0}^{\infty}A_{\varepsilon}(r)L(r)\,\xi_{r,\varepsilon}(w)\,dr,\label{eq:T1}\\
 T_{2}^{\varepsilon}(w\mid Y)
 &=\bigl(\rho_{\varepsilon}*B_{\varepsilon}L\bigr)(w),\label{eq:T2}\\
 T_{3}^{\varepsilon}(w\mid Y)
 &=-\partial_{w}\bigl(\rho_{\varepsilon}*C_{\varepsilon}L\bigr)(w),\label{eq:T3}
\end{align}
where $T_j^{\varepsilon}$ use the bulk-extrapolated coefficients
integrated over all $[0,\infty)$, and the boundary-layer residual
\begin{equation}\label{eq:R-layer}
 R_{\varepsilon}^{\mathrm{layer}}(w\mid Y)
 =\tfrac12\!\int_{0}^{c\varepsilon}\!\!\bigl[\Delta_{A}(r)\,\xi_{r,\varepsilon}(w)
 +\Delta_{B}(r)\,\rho_{\varepsilon}(r-w)\bigr]L(r)\,dr
\end{equation}
captures the deviation from the bulk identity on $r\in[0,c\varepsilon]$,
with
$\Delta_{A}(r)=(\Var(V_r)+\Cov(U_{r},W_{r}))\E_{X}[f'''(x+U_{r})]
+(\Cov(U_{r},V_{r})^{2}-\tfrac14)\E_{X}[f^{(5)}(x+U_{r})]$,
$\Delta_{B}(r)=(2\Cov(U_{r},V_{r})-1)\E_{X}[f'''(x+U_{r})]$. We write
$\Var(V_r)$ here rather than $\sigma_V^{2}$ to emphasise that the
boundary-layer factor must be read as the pointwise variance; the
two coincide in the bulk by~\eqref{eq:bulk-cov}. Both
$\Delta_A$, $\Delta_B$ are supported in $[0,c\varepsilon]$ by
\eqref{eq:bulk-cov}; the boundary covariance
$\Cov(U_0,V_0)=0$ (\cref{lem:bdy-cov-zero}) is used in
\eqref{eq:Kbdy} to kill the $\xi_{0,\varepsilon}$ contribution at
$r=0$.

The boundary-layer residual $R_\varepsilon^{\mathrm{layer}}$
captures the deviation from the bulk Stein identity on the
$O(\varepsilon)$-neighbourhood of the origin. It vanishes in the
$L^{2}$ topology with an explicit $\varepsilon^{1/2}$ rate; this is
the input the interior convergence theorem needs.

\begin{lemma}[Boundary-layer residual vanishing]\label{lem:bdy-layer-vanish}
For every $f\in\mathscr S(\R)$,
\[
 R_{\varepsilon}^{\mathrm{layer}}\;\xrightarrow[\varepsilon\to 0]{}\;0
 \quad\text{in }L^{2}\bigl(\Omega_{Y};\,L^{2}(\R,dw)\bigr),
\]
and hence in $L^{2}(\Omega_{Y};H^{-1}(\R))$ via the continuous
embedding $L^{2}\hookrightarrow H^{-1}$. Moreover the bound
\begin{equation}\label{eq:R-layer-bound}
 \|R_{\varepsilon}^{\mathrm{layer}}\|_{L^{2}(\R)}
 \;\le\;C(f,\rho)\,\varepsilon^{1/2}\,\|L\|_{\infty}
 \quad\text{a.s.\ in }Y
\end{equation}
holds with $C(f,\rho)$ depending only on
$\|f'''\|_{\infty},\|f^{(5)}\|_{\infty}$, $c$, $\|\rho\|_{\infty}$,
and $\int\rho^{2}$.
\end{lemma}

\begin{proof}
\emph{Step 1: closed form for $(\sigma_{\varepsilon}^{2})''$.}
Change of variables in
$\sigma_{\varepsilon}^{2}(r)=\iint\rho_{\varepsilon}(r-z_{1})\rho_{\varepsilon}(r-z_{2})\Cov(X_{z_{1}},X_{z_{2}})\,dz_{1}dz_{2}$
(using two-sided scenery covariance) yields
$\sigma_{\varepsilon}^{2}(r)=(\rho_{\varepsilon}*|\cdot|)(r)-K_{\rho}\varepsilon$,
so $(\sigma_{\varepsilon}^{2})'(r)=(\rho_{\varepsilon}*\sgn)(r)$ and
$(\sigma_{\varepsilon}^{2})''(r)=2\rho_{\varepsilon}(r)$, using
$(\sgn)'=2\delta_0$ distributionally.

\emph{Step 2: layer factor bounds.}
Differentiating $\Var(U_r)=\sigma_\varepsilon^{2}(r)$ twice gives
$2(\Var(V_r)+\Cov(U_r,W_r))=(\sigma_\varepsilon^{2})''(r)$, so
$|\Var(V_r)+\Cov(U_r,W_r)|=\rho_\varepsilon(r)\le\|\rho\|_\infty/\varepsilon$
for $r\in[0,c\varepsilon]$, vanishing in the bulk. Similarly
$2\Cov(U_r,V_r)=(\rho_\varepsilon*\sgn)(r)\in[-1,1]$, equal to $0$
at $r=0$ (\cref{lem:bdy-cov-zero}); for $r\in[0,c\varepsilon]$,
$\Cov(U_r,V_r)\in[0,1/2]$, so
$|\Cov(U_r,V_r)^{2}-\tfrac14|\le\tfrac14$ and
$|2\Cov(U_r,V_r)-1|\le 1$. With $f^{(k)}$ bounded for
$f\in\mathscr S(\R)$,
\(
 |\Delta_A(r)|\le C_f\rho_\varepsilon(r)+\tfrac14\|f^{(5)}\|_\infty,
\)
\(
 |\Delta_B(r)|\le\|f'''\|_\infty,
\)
where $C_f=\|f'''\|_\infty$.

\emph{Step 3: $L^{2}(\R)$ bound.}
By \cref{lem:moll-H}(a), $\|\xi_{r,\varepsilon}\|_{L^{2}}\le(c+C_\rho)^{1/2}\varepsilon^{1/2}$
for $r\in[0,c\varepsilon]$. Also
$\|\rho_\varepsilon\|_{L^{2}}=(\int\rho^{2})^{1/2}\varepsilon^{-1/2}$.
Triangle inequality in $L^{2}(\R,dw)$ applied to~\eqref{eq:R-layer}:
\begin{align*}
 \|R_{\varepsilon}^{\mathrm{layer}}\|_{L^{2}(\R)}
 &\le\tfrac12(c+C_\rho)^{1/2}\varepsilon^{1/2}\|L\|_\infty\!\int_0^{c\varepsilon}|\Delta_A(r)|\,dr\\
 &\quad+\tfrac12(\textstyle\int\rho^{2})^{1/2}\varepsilon^{-1/2}\|L\|_\infty\!\int_0^{c\varepsilon}|\Delta_B(r)|\,dr.
\end{align*}
With $\int_0^{c\varepsilon}\rho_\varepsilon\le 1$ and
$\int_0^{c\varepsilon}1=c\varepsilon$,
$\int_0^{c\varepsilon}|\Delta_A|\le C_f+\tfrac14\|f^{(5)}\|_\infty c\varepsilon$,
$\int_0^{c\varepsilon}|\Delta_B|\le\|f'''\|_\infty c\varepsilon$.
Substituting yields~\eqref{eq:R-layer-bound}.

\emph{Step 4: assembly and conclusion.}
Squaring \eqref{eq:R-layer-bound} and taking $\E_Y$,
$\E_Y\|R_\varepsilon^{\mathrm{layer}}\|_{L^{2}}^{2}
\le C(f,\rho)^{2}\varepsilon\,\E_Y[\|L\|_\infty^{2}]$.
$\E_Y[\|L\|_\infty^{2}]<\infty$ via
$L_t^{r}(|Y|)=L_t^{r}(Y)+L_t^{-r}(Y)$ and the finite-moment estimate
$\E[\sup_x L_t^{x}(Y)^{p}]<\infty$ (\cite[Thm.\ 9.4.1]{MarcusRosen}).
Hence convergence to $0$ at rate $O(\varepsilon)$, and the
$L^{2}\hookrightarrow H^{-1}$ embedding gives the second statement.
\end{proof}

\subsection*{Interior convergence in $H^{-1}$.}

The qualitative $H^{-1}$ convergence of the interior kernel is
unconditional in $f\in\mathscr S(\R)$. A quantitative rate, given
at the end of the subsection, requires an additional
Marcus--Rosen-type moment estimate on the reflected local time,
isolated in the following assumption.

For a function $h\colon\R_+\to\R$ supported in $[0,M_t]$, we
interpret $\|h\|_{H^{1/2-\eta}([0,M_t])}$ as the
$H^{1/2-\eta}(\R)$-norm of the zero-extension of $h$ to the whole
real line. This is the norm in which the Plancherel-based mollifier
estimate $\|\rho_\varepsilon*g-g\|_{L^{2}}\le C_\rho\varepsilon^{s}\|g\|_{H^{s}}$
operates downstream in \cref{prop:interior-rate}.

\begin{assumption}[Local-time Sobolev moments]\label{assu:lt-moment}
For every $\eta\in(0,1/2)$ and $p\ge 2$,
\(
 \E_{Y}\bigl[\|L_{t}^{\cdot}(|Y|)\|_{H^{1/2-\eta}([0,M_{t}])}^{p}\bigr]<\infty,
\)
with the random-interval norm read as the whole-line
$H^{1/2-\eta}(\R)$-norm of the zero-extension as above.
\end{assumption}

\Cref{assu:lt-moment} is consistent with standard Marcus--Rosen
local-time moment estimates (cf.~\cite{MarcusRosen}); we isolate it
as an assumption because the precise Sobolev-moment formulation,
read against the whole-line zero-extension above, is not needed for
qualitative convergence. The qualitative convergence
\cref{thm:interior,thm:boundary,cor:combined} does not invoke
\cref{assu:lt-moment}; only the rate \cref{prop:interior-rate} does.

Define the \emph{bulk-extrapolated interior kernel} and the
\emph{actual interior kernel}:
\[
 K_{\varepsilon}^{\mathrm{int}}\;\coloneqq\;\tfrac12(T_{1}^{\varepsilon}+T_{2}^{\varepsilon}+T_{3}^{\varepsilon}),
 \qquad
 K_{\varepsilon}^{\mathrm{int,exact}}\;\coloneqq\;K_{\varepsilon}^{\mathrm{int}}+R_{\varepsilon}^{\mathrm{layer}}.
\]

\begin{theorem}[Interior first-chaos $H^{-1}$ convergence]
\label{thm:interior}
Let $f\in\mathscr S(\R)$. Then both $K_{\varepsilon}^{\mathrm{int}}$
and $K_{\varepsilon}^{\mathrm{int,exact}}$ converge to the same limit
\[
 K_{\varepsilon}^{\mathrm{int}},\,K_{\varepsilon}^{\mathrm{int,exact}}
 \;\xrightarrow[\varepsilon\to 0]{}\;K^{\mathrm{int}}
 \quad\text{in }L^{2}\bigl(\Omega_{Y};\,H^{-1}(\R,dw)\bigr),
\]
with
\begin{equation}\label{eq:Kint-limit}
 K^{\mathrm{int}}(w\mid Y)
 =\tfrac12\bigl[T_{1}^{*}(w\mid Y)+B(w)L(w)-\partial_{w}\bigl(C(w)L(w)\bigr)\bigr],
\end{equation}
where $T_{1}^{*}(w\mid Y)=\mathbf 1_{w\ge 0}\int_{w}^{M_{t}}A(r)L(r)\,dr$,
$L=L_{t}^{\cdot}(|Y|)$, $M_{t}=\sup_{[0,t]}|Y_{s}|$, and $A,B,C$ are
the heat-semigroup limits of $A_\varepsilon,B_\varepsilon,C_\varepsilon$.
\end{theorem}

\Cref{thm:interior} does not use the nondegeneracy
hypothesis~\eqref{eq:nondegen}; convergence holds for every
$f\in\mathscr S(\R)$. The role of~\eqref{eq:nondegen} is to certify
that the present $H^{-1}$-convergent sector is the same sector whose
$L^{2}(\Omega)$-variance diverges in~\cref{thm:var-divergence}:
under~\eqref{eq:nondegen}, $C\not\equiv 0$ on $(0,\infty)$, so the
$T_3^{*}=-\partial_w(CL)$ summand is non-trivial, and the divergence
in~\eqref{eq:var-diverges} is the $L^{2}(\R)$-norm failure of the
same summand. Without~\eqref{eq:nondegen} the
limit~\eqref{eq:Kint-limit} still exists in $H^{-1}$; it is simply
not witnessed by an $L^{2}$ obstruction.

\begin{proof}
By the triangle inequality,
$K_\varepsilon^{\mathrm{int,exact}}-K^{\mathrm{int}}
=(K_\varepsilon^{\mathrm{int}}-K^{\mathrm{int}})+R_\varepsilon^{\mathrm{layer}}$,
and $R_\varepsilon^{\mathrm{layer}}\to 0$ in
$L^{2}(\Omega_Y;L^{2}(\R))\hookrightarrow L^{2}(\Omega_Y;H^{-1}(\R))$
by \cref{lem:bdy-layer-vanish}. It suffices to prove
$K_\varepsilon^{\mathrm{int}}\to K^{\mathrm{int}}$ in
$L^{2}(\Omega_Y;H^{-1}(\R))$, using
$\|\cdot\|_{H^{-1}}\le\|\cdot\|_{L^{2}}$.

Decompose
$K_\varepsilon^{\mathrm{int}}-K^{\mathrm{int}}=\tfrac12\sum_{j=1}^{3}\Delta_j^{\varepsilon}$
with $\Delta_j^{\varepsilon}=T_j^{\varepsilon}-T_j^{*}$.

\emph{(a) $T_1^{\varepsilon}\to T_1^{*}$ in $L^{2}(\Omega_Y;L^{2}(\R))$.}
Split $\Delta_1^{\varepsilon}(w)=I_\varepsilon(w)+J_\varepsilon(w)$
with
$I_\varepsilon=\int_0^\infty[A_\varepsilon(r)-A(r)]L(r)\xi_{r,\varepsilon}(w)\,dr$,
$J_\varepsilon=\int_0^\infty A(r)L(r)[\xi_{r,\varepsilon}(w)-\sigma_r(w)]\,dr$.
$L$ has compact support in $[0,M_t]$. Heat-semigroup
time-Lipschitz~(\cref{lem:heat-rate} below) gives
$A_\varepsilon\to A$ uniformly on $[0,M_t]$ a.s.\ in $Y$;
\cref{lem:moll-H}(a) gives
$\|\xi_{r,\varepsilon}\|_{L^{2}}\le(M_t+C_\rho)^{1/2}$. Minkowski in
$w$:
$\|I_\varepsilon\|_{L^{2}}\le\|A_\varepsilon-A\|_{L^{\infty}([0,M_t])}\cdot t\cdot(M_t+C_\rho)^{1/2}\to 0$
a.s.\ in $Y$. For $J_\varepsilon$, \cref{lem:moll-H}(b) gives
$\|\xi_{r,\varepsilon}-\sigma_r\|_{L^{2}}\le C'_\rho\varepsilon^{1/2}$,
so $\|J_\varepsilon\|_{L^{2}}\le\|A\|_\infty\cdot t\cdot C'_\rho\varepsilon^{1/2}\to 0$
a.s.\ in $Y$. Both bounds are dominated by
$C(f,\rho)t(1+M_t^{1/2})\in L^{2}(\Omega_Y)$ (classical
$\E_Y[M_t^{p}]<\infty$); dominated convergence concludes.

\emph{(b) $T_2^{\varepsilon}\to BL$ in $L^{2}(\Omega_Y;L^{2}(\R))$.}
$\Delta_2^{\varepsilon}=\rho_\varepsilon*((B_\varepsilon-B)L)+(\rho_\varepsilon*BL-BL)$.
Young: $\|\rho_\varepsilon*((B_\varepsilon-B)L)\|_{L^{2}}\le\|B_\varepsilon-B\|_\infty\|L\|_{L^{2}}\to 0$
a.s.\ in $Y$. Standard $L^{2}$-mollifier convergence:
$\|\rho_\varepsilon*BL-BL\|_{L^{2}}\to 0$ a.s. Both terms dominated
by $\|B\|_\infty\|L\|_{L^{2}}\in L^{2}(\Omega_Y)$; dominated
convergence concludes.

\emph{(c) $T_3^{\varepsilon}\to-\partial_w(CL)$ in
$L^{2}(\Omega_Y;H^{-1}(\R))$.}
$\partial_w\colon L^{2}\to H^{-1}$ has norm $\le 1$, so
$\|\Delta_3^{\varepsilon}\|_{H^{-1}}\le\|\rho_\varepsilon*C_\varepsilon L-CL\|_{L^{2}}\to 0$
a.s.\ in $Y$ by the same argument as (b) applied to $g=CL$;
domination by $\|C\|_\infty\|L\|_{L^{2}}$ finishes.
\end{proof}

\begin{lemma}[Heat-semigroup time-Lipschitz]\label{lem:heat-rate}
For $h\in\mathscr S(\R)$ and $\eta_1,\eta_2\ge 0$,
$\|P_{\eta_1}h-P_{\eta_2}h\|_{L^{\infty}(\R)}\le\tfrac12|\eta_1-\eta_2|\,\|h''\|_{L^{\infty}}$.
In particular
$|C_\varepsilon(r)-C(r)|\le C_f|\sigma_\varepsilon^{2}(r)-|r||\le C_f\varepsilon$
uniformly in $|r|\ge c\varepsilon$, and
$|C_\varepsilon(0)-f'(x)|\le C_f\varepsilon$.
\end{lemma}

Refining the three $\Delta_j^{\varepsilon}$ estimates of
\cref{thm:interior} with \cref{assu:lt-moment} promotes the
qualitative limit to a quantitative rate.

\begin{proposition}[Quantitative interior rate, conditional]\label{prop:interior-rate}
Assume \cref{assu:lt-moment}. Then for every $\eta\in(0,1/2)$,
\begin{equation}\label{eq:int-rate}
 \bigl\|K_{\varepsilon}^{\mathrm{int,exact}}-K^{\mathrm{int}}\bigr\|_{L^{2}(\Omega_{Y};H^{-1}(\R))}
 \;\le\;C(f,t,\rho,\eta)\,\varepsilon^{1/2-\eta}.
\end{equation}
The same bound holds with $K_\varepsilon^{\mathrm{int}}$ in place
of $K_\varepsilon^{\mathrm{int,exact}}$. The rate is conditional on
\cref{assu:lt-moment}; the qualitative convergence
(\cref{thm:interior}) is not.
\end{proposition}

\begin{proof}
The residual is $O(\varepsilon^{1/2})$ in $L^{2}(\Omega_Y;L^{2}(\R))$
by~\eqref{eq:R-layer-bound} and $\E_Y[\|L\|_\infty^{2}]<\infty$;
hence in $L^{2}(\Omega_Y;H^{-1}(\R))$ at the same rate. We refine the
three $\Delta_j^{\varepsilon}$ estimates of \cref{thm:interior}.

\emph{(a) Rate for $\Delta_1^{\varepsilon}$.}
\Cref{lem:heat-rate} gives $\|A_\varepsilon-A\|_{L^{\infty}([0,M_t])}\le C_f\varepsilon$;
\cref{lem:moll-H}(b) gives
$\|\xi_{r,\varepsilon}-\sigma_r\|_{L^{2}}\le C'_\rho\varepsilon^{1/2}$.
Combining,
$\|\Delta_1^{\varepsilon}\|_{L^{2}}^{2}\le C(f,\rho)^{2}\,t^{2}(1+M_t)\varepsilon$
a.s., so $\E_Y\|\Delta_1^{\varepsilon}\|_{L^{2}}^{2}\le C(f,\rho)^{2}t^{2}(1+\E_Y M_t)\varepsilon$.

\emph{(b) Rate for $\Delta_2^{\varepsilon}$.}
$\Delta_2^{\varepsilon}=\rho_\varepsilon*((B_\varepsilon-B)L)+(\rho_\varepsilon*BL-BL)$.
By Young, $\|\rho_\varepsilon*((B_\varepsilon-B)L)\|_{L^{2}}\le C_f\varepsilon\|L\|_{L^{2}}$.
For the second summand, the mollifier $H^{s}$-rate
$\|\rho_\varepsilon*g-g\|_{L^{2}}\le C_\rho\varepsilon^{s}\|g\|_{H^{s}}$
for $s\in(0,1]$ (Plancherel) applied with $g=BL$, $s=1/2-\eta$, plus
$\|BL\|_{H^{1/2-\eta}}\le C_f\|L\|_{H^{1/2-\eta}}$ ($B$ smooth
bounded) and \cref{assu:lt-moment}, gives
$\E_Y\|\Delta_2^{\varepsilon}\|_{L^{2}}^{2}\le C(f,\rho)^{2}\varepsilon^{1-2\eta}$.

\emph{(c) Rate for $\Delta_3^{\varepsilon}$ in $H^{-1}$.}
Same argument as (b) for $g=CL$, $s=1/2-\eta$, then
$\|\partial_w\|_{L^{2}\to H^{-1}}\le 1$:
$\E_Y\|\Delta_3^{\varepsilon}\|_{H^{-1}}^{2}\le C(f,\rho)^{2}\varepsilon^{1-2\eta}$.

$\|\cdot\|_{H^{-1}}\le\|\cdot\|_{L^{2}}$ on (a) and (b); summing
yields \eqref{eq:int-rate}.
\end{proof}

The rate is consistent with the sharp Brownian local-time Sobolev
threshold at $1/2^{-}$: $L$ has spatial Sobolev regularity exactly
up to $H^{1/2-\eta}$, and the mollifier rate
$\|\rho_\varepsilon*g-g\|_{L^{2}}\sim\varepsilon^{s}$ for
$g\in H^{s}$ is sharp at the regularity index of $g$. The
$\eta$-loss is intrinsic to Brownian local-time regularity, not an
artefact of the proof. No endpoint rate is claimed. The rate depends
on the moment input~\cref{assu:lt-moment}; the qualitative
convergence does not.

\subsection*{Boundary contribution and the combined closure.}

The $T_1$ and $T_2$ pieces of $K^{\mathrm{int}}$ are in $L^{2}(\R)$
on the qualitative limit, while the $T_3$ piece $-\partial_w(CL)$
is generally only in $H^{-1}(\R)$; this is the failure identified
in~\cref{thm:var-divergence}. The boundary kernel
$K_\varepsilon^{\mathrm{bdy}}=L_t^{0}(Y)C_\varepsilon(0)\rho_\varepsilon$
collapses to the singular distribution $L_t^{0}(Y)f'(x)\delta_0$.
Since $\delta_0\notin H^{-1/2}$ but $\delta_0\in H^{-1/2-\delta}$
for every $\delta>0$, the boundary convergence is most naturally
stated one Sobolev step weaker than $H^{-1}$; the inclusion
$H^{-1/2-\delta}\subset H^{-1}$ for $\delta\in(0,1/2)$ then assembles
the combined statement.

\begin{theorem}[Boundary first-chaos $H^{-1/2-\delta}$ convergence]
\label{thm:boundary}
For every $f\in\mathscr S(\R)$ and every $\delta\in(0,1/2)$,
\[
 K_{\varepsilon}^{\mathrm{bdy}}\;\xrightarrow[\varepsilon\to 0]{}\;L_{t}^{0}(Y)\,f'(x)\,\delta_{0}
 \quad\text{in }L^{2}\bigl(\Omega_{Y};\,H^{-1/2-\delta}(\R)\bigr),
\]
with rate
$\bigl\|K_{\varepsilon}^{\mathrm{bdy}}-L_{t}^{0}(Y)f'(x)\delta_{0}\bigr\|_{L^{2}(\Omega_{Y};H^{-1/2-\delta}(\R))}
\le C_{f,\rho}\,\varepsilon^{\delta\wedge 1/2}$. The range
$\delta\in(0,1/2)$ matches the Fourier-estimate range used in the
proof and is the range needed by~\cref{cor:combined} for the
embedding $H^{-1/2-\delta}\subset H^{-1}$. No $\delta_0'$
contribution arises (\cref{lem:bdy-cov-zero,rem:bdy-coeff}).
\end{theorem}

\begin{proof}
$K_\varepsilon^{\mathrm{bdy}}=L_t^{0}(Y)C_\varepsilon(0)\rho_\varepsilon$.
Split
$K_\varepsilon^{\mathrm{bdy}}-L_t^{0}(Y)f'(x)\delta_0
=L_t^{0}(Y)(C_\varepsilon(0)-f'(x))\rho_\varepsilon+L_t^{0}(Y)f'(x)(\rho_\varepsilon-\delta_0)$.

\emph{Step 1 (coefficient rate):} $C_\varepsilon(0)=(P_{\sigma_\varepsilon^{2}(0)}f')(x)$;
$\sigma_\varepsilon^{2}(0)=O(\varepsilon)$ by Step~1 of
\cref{lem:bdy-layer-vanish}. \Cref{lem:heat-rate} gives
$|C_\varepsilon(0)-f'(x)|\le C_f\varepsilon$.

\emph{Step 2 (distribution rate):} Plancherel and
$\hat\rho_\varepsilon(k)=\hat\rho(\varepsilon k)$,
$|\hat\rho(\varepsilon k)-1|\le\min(2,c_\rho\varepsilon^{2}k^{2})$,
together with the split at $|k|=1/\varepsilon$:
$\|\rho_\varepsilon-\delta_0\|_{H^{-s}}\le C\varepsilon^{s-1/2}$ for
$s\in(1/2,5/2)$. Applied at $s=1/2+\delta$:
$\|\rho_\varepsilon-\delta_0\|_{H^{-1/2-\delta}}\le C\varepsilon^{\delta}$.

\emph{Step 3 (assemble):} Triangle inequality and
$\E[L_t^{0}(Y)^{2}]<\infty$ (\cite[Thm.\ 6.1.1]{MarcusRosen}) give the
claimed bound, with the slower component
$\varepsilon^{\delta\wedge 1/2}$.
\end{proof}

The Sobolev exponent in~\cref{thm:boundary} cannot be improved at
the boundary: $\delta_0\in H^{s}(\R)$ if and only if $s<-1/2$, so
$\delta_0\notin H^{-1/2}$ but $\delta_0\in H^{-1/2-\delta}$ for
every $\delta>0$. We work at $\delta\in(0,1/2)$, which both matches
the Fourier-estimate range in the proof and gives the inclusion
$H^{-1/2-\delta}\subset H^{-1}$ needed to assemble the combined
$H^{-1}$ closure of the full first-$X$-chaos Malliavin kernel.

\begin{corollary}[Combined first-chaos $H^{-1}$ convergence]
\label{cor:combined}
For every $f\in\mathscr S(\R)$ and every $\delta\in(0,1/2)$, the
full first-$X$-chaos Malliavin kernel
\(K_{\varepsilon}=K_{\varepsilon}^{\mathrm{bdy}}+K_{\varepsilon}^{\mathrm{int,exact}}\)
satisfies
\[
 K_{\varepsilon}\;\xrightarrow[\varepsilon\to 0]{}\;K
 \;\coloneqq\;L_{t}^{0}(Y)f'(x)\delta_{0}+K^{\mathrm{int}}
 \quad\text{in }L^{2}\bigl(\Omega_{Y};\,H^{-1}(\R,dw)\bigr).
\]
Under \cref{assu:lt-moment}, the combined rate is
$\varepsilon^{(1/2-\eta)\wedge\delta}$ for every $\eta\in(0,1/2)$ and
$\delta\in(0,1/2)$; choosing $\delta=1/2-\eta$ optimises to
$\varepsilon^{1/2-\eta}$.
\end{corollary}

\begin{proof}
Combine \cref{thm:interior,thm:boundary}. The embedding
$H^{-1/2-\delta}\hookrightarrow H^{-1}$ for $\delta\in(0,1/2)$ is
standard; the interior $L^{2}\hookrightarrow H^{-1}$ embedding is
direct.
\end{proof}

\begin{corollary}[First-chaos dual closure of $\Pir_1^{X}C_{\IBM,\varepsilon}$]
\label{cor:proj1-closure}
At every $\varepsilon>0$, $\Pir_1^{X}C_{\IBM,\varepsilon}(f;t)$ is
an honest first-chaos random variable with conditional Malliavin
kernel $K_\varepsilon\in L^{2}(\Omega_{Y};L^{2}(\R))$;
\cref{cor:combined} gives $K_\varepsilon\to K$ in
$L^{2}(\Omega_{Y};H^{-1}(\R))=\Energy_{\IBM,1}^{-}$, with limit
kernel $K$ generically only $H^{-1}$ (not $L^{2}$), so the limit
object is a generalized $Y$-conditioned first-chaos functional
(\cref{def:first-chaos-block}) rather than an $L^{2}(\Omega_X)$
random variable. Convergence is the duality-pairing convergence:
for every $\Phi\in\Energy_{\IBM,1}^{+}=L^{2}(\Omega_Y;H^{1}(\R))$,
\[
 \bigl\langle\Pir_{1}^{X}C_{\IBM,\varepsilon}(f;t),\Phi\bigr\rangle_{1,*}
 \;\xrightarrow[\varepsilon\to 0]{}\;
 \E_{Y}\!\bigl[\langle K(\cdot\mid Y),\Phi(\cdot\mid Y)\rangle_{H^{-1},H^{1}}\bigr],
\]
including the special case of deterministic $\phi\in H^{1}(\R)$.
Under \cref{assu:lt-moment}, the rate is $\varepsilon^{1/2-\eta}$.
\end{corollary}

\Cref{cor:combined,cor:proj1-closure} close the
$L^{2}(\Omega)$-divergent first scenery-chaos sector of
$C_{\IBM,\varepsilon}(f;t)$ in the first-$X$-chaos block of the
operator-derivative dual. They do not prove:
\begin{itemize}
\item that $C_\IBM(f;t)$ exists in any global cross-sector dual; the
$(m,n)$-bidegree projections with $m\ge 2$ are not addressed here
(\cref{flag:summability});
\item that the full unmollified IBM It\^o
identity~\eqref{eq:flagship-intro} survives $\varepsilon\to 0$ in
any topology;
\item that the Brownian-time semigroup correction
$u(t,x)-f(x)$ arises from the first-chaos sector. It cannot:
first chaos has mean zero, so $\E[\Pir_1^{X}C_{\IBM}(f;t)]=0$. The
semigroup correction lives in $\Pir_{0,0}$ (\cref{thm:pure-clock}).
\end{itemize}
Filtration-level obstructions (Tsirelson-type non-Brownian noises)
are an independent direction and play no role here. The next
section places this closure inside the joint Wiener bidegree
lattice and shows how the same construction, at fixed $\varepsilon$,
recovers the Allouba--Zheng semigroup correction in the pure-clock
sector.

%==============================================================
\section{The bidegree atlas at fixed $\varepsilon$}\label{sec:atlas}

This section is supporting structure rather than a fourth main
result: it provides a fixed-$\varepsilon$ structural map of
$C_{\IBM,\varepsilon}(f;t)$ across the joint Wiener chaos lattice
$\Hilb_m^X\otimes\Hilb_n^Y$, locating the first-chaos dual closure
of~\cref{sec:dual} inside the lattice and exhibiting the sector
that carries the Brownian-time semigroup correction
of~\cref{thm:moll-consistency}. At fixed
$\varepsilon$, $C_{\IBM,\varepsilon}\in L^{2}(\Omega_X\times\Omega_Y)$
and the joint chaos expansion
\begin{equation}\label{eq:atlas-conv}
 C_{\IBM,\varepsilon}(f;t)
 \;=\;\sum_{m,n\ge 0}\Pir_{m,n}\,C_{\IBM,\varepsilon}(f;t)
\end{equation}
converges in $L^{2}(\Omega)$ by Wiener--It\^o
(\cref{thm:atlas} below). All statements that follow are at fixed
$\varepsilon>0$; cross-sector $\varepsilon\to 0$
summability is the open problem \cref{flag:summability}. Three
sectors carry the structural content: pure-clock $(0,\bullet)$,
first-scenery-chaos $(1,\bullet)$, and the $(2,0)$ scalar
energy-contraction channel.

\begin{theorem}[Bidegree decomposition at every $\varepsilon$]
\label{thm:atlas}
For every $f\in\mathscr S(\R)$, $t\ge 0$, $\varepsilon>0$,
$C_{\IBM,\varepsilon}(f;t)\in L^{2}(\Omega_X\times\Omega_Y)$ and
\eqref{eq:atlas-conv} converges in $L^{2}(\Omega_X\times\Omega_Y)$.
The decomposition applies summand-by-summand to the two terms
$g_\varepsilon'(0)L_t^0(Y)$ and
$\tfrac12\!\int_0^t g_\varepsilon''(|Y_s|)\,ds$ of
$C_{\IBM,\varepsilon}$.
\end{theorem}

\begin{proof}
At fixed $\varepsilon>0$, each of $g_\varepsilon'(0)$,
$g_\varepsilon''(r)$ has uniform $L^{2}(\Omega_X)$ bounds in $r$
(smooth bounded $f^{(k)}(x+U_r)$ for $k\le 5$, Gaussian factors
$V_r,V_r^{2},W_r$ in $L^{p}$ at fixed $\varepsilon$);
$L_t^{0}(Y)\in L^{p}(\Omega_Y)$ for every $p$; Minkowski preserves
$L^{2}$ under $\int_0^t$. Hence $C_{\IBM,\varepsilon}(f;t)\in L^{2}(\Omega)$,
and \eqref{eq:atlas-conv} is the orthogonal expansion in
\eqref{eq:joint-chaos}, convergent by Wiener--It\^o.
\end{proof}

\subsection*{The pure-clock sector $(0,\bullet)$.}

The pure-clock sector is $X$-deterministic; it carries the
expectation-level content of the correction. Define the full
$X$-expectation of the scalar integrand:
\begin{equation}\label{eq:psi-def}
 \psi_\varepsilon(r)\;\coloneqq\;\E_X[g_\varepsilon''(r)]
 \;=\;\E_X[f''(x+U_r)V_r^{2}+f'(x+U_r)W_r].
\end{equation}
Stein on $(U_r,V_r)$ and $(U_r,W_r)$ gives
\begin{equation}\label{eq:psi-decomp}
 \psi_\varepsilon(r)
 \;=\;\Cov(U_r,V_r)^{2}\,\E_X[f^{(4)}(x+U_r)]
   +\bigl[\Var(V_r)+\Cov(U_r,W_r)\bigr]\,\E_X[f''(x+U_r)],
\end{equation}
where we write $\Var(V_r)$ in place of $\sigma_V^{2}$ to keep the
boundary-layer factor pointwise; the two coincide in the bulk
by~\eqref{eq:bulk-cov}. By~\eqref{eq:bulk-cov},
$\Cov(U_r,V_r)^{2}=\tfrac14$ and $\Var(V_r)+\Cov(U_r,W_r)=0$ for
$r\ge c\varepsilon$. Setting
\begin{equation}\label{eq:phi-int-def}
 \phi_\varepsilon^{\mathrm{int}}(r)\coloneqq\tfrac14\E_X[f^{(4)}(x+U_r)],
\qquad
 b_\varepsilon(r)\coloneqq\psi_\varepsilon(r)-\phi_\varepsilon^{\mathrm{int}}(r),
\end{equation}
the boundary residual $b_\varepsilon$ is smooth, supported in
$[0,c\varepsilon]$, and satisfies (by~\eqref{eq:psi-decomp}
and bracket bounds on the layer)
\begin{equation}\label{eq:b-eps-bounds}
 \sup_r|b_\varepsilon(r)|\le C_f\varepsilon^{-1},\qquad
 \int_0^\infty|b_\varepsilon(r)|\,dr\le C_f,
\end{equation}
with $C_f$ depending on $\|f^{(2)}\|_\infty,\|f^{(4)}\|_\infty$ and
$\rho$ only. The bound~\eqref{eq:b-eps-bounds} follows by the same
bracket argument used in \cref{lem:bdy-layer-vanish}: the
boundary-layer factors $|\Var(V_r)+\Cov(U_r,W_r)|$ and
$|\Cov(U_r,V_r)^{2}-\tfrac14|$ are bounded by
$\rho_\varepsilon(r)\le\|\rho\|_\infty/\varepsilon$ and $\tfrac14$
respectively for $r\in[0,c\varepsilon]$, both vanishing on
$\{r\ge c\varepsilon\}$.

\begin{lemma}[Boundary Tanaka piece vanishes after $\E_X$]
\label{lem:bdy-zero}
$\E_X[g_\varepsilon'(0)]=0$.
\end{lemma}

\begin{proof}
$\E_X[g_\varepsilon'(0)]=\E_X[f'(x+U_0)V_0]=\Cov(U_0,V_0)\E_X[f''(x+U_0)]=0$
by the scalar Stein identity and the boundary covariance
$\Cov(U_0,V_0)=0$ (\cref{lem:bdy-cov-zero}).
\end{proof}

\begin{theorem}[Pure-clock sector and expectation limit]
\label{thm:pure-clock}
For every $f\in\mathscr S(\R)$, $t\ge 0$, $\varepsilon>0$:
\begin{enumerate}
\item (Scalar identity)
\[
 \Pir_0^X\,C_{\IBM,\varepsilon}(f;t)
 \;=\;F_\varepsilon(Y;t)
 \;\coloneqq\;\tfrac12\!\int_0^t\!\psi_\varepsilon(|Y_s|)\,ds,
\]
with $\psi_\varepsilon=\phi_\varepsilon^{\mathrm{int}}+b_\varepsilon$
as in \eqref{eq:phi-int-def}.
\item (Bidegree slices) For every $n\ge 0$,
$\Pir_{0,n}\,C_{\IBM,\varepsilon}(f;t)=\Pir_n^Y F_\varepsilon(Y;t)$.
\item (Expectation limit, referring to \cref{thm:moll-consistency})
\[
 \Pir_{0,0}\,C_{\IBM,\varepsilon}(f;t)
 \;=\;\E\!\bigl[C_{\IBM,\varepsilon}(f;t)\bigr]
 \;\xrightarrow{\;\varepsilon\to 0\;}\;u(t,x)-f(x).
\]
The limit decomposes as
\[
 u(t,x)-f(x)
 \;=\;\underbrace{\lim_{\varepsilon\to 0}\tfrac12\E_Y\!\int_0^t\phi_\varepsilon^{\mathrm{int}}(|Y_s|)\,ds}_{\text{interior bi-Laplacian shadow}}
   \;+\;\underbrace{\lim_{\varepsilon\to 0}\tfrac12\E_Y\!\int_0^t b_\varepsilon(|Y_s|)\,ds}_{\text{boundary-layer source}};
\]
the boundary-layer piece converges in expectation to
$\int_0^t\partial_x^{2}f(x)/(2\sqrt{2\pi s})\,ds$. We identify this
limit only at the level of expectations; see \cref{flag:bdy-layer}.
\item (Clock-fluctuation blocks) For $n\ge 1$,
$\Pir_{0,n}=\Pir_n^Y F_\varepsilon$ is centred in $\Hilb_n^Y$. In
particular, by Stroock--Clark--Ocone~\cite{Nualart},
\[
 \Pir_{0,1}\,C_{\IBM,\varepsilon}(f;t)
 \;=\;I_1^Y\!\bigl(\E_Y[D^Y_\cdot F_\varepsilon(Y;t)]\bigr),
\]
with $D^Y_u|Y_s|=\sgn(Y_s)\mathbf 1_{u\le s}$ on the standard
Bouleau--Hirsch domain.
\end{enumerate}
\end{theorem}

\begin{proof}
(1) $\Pir_0^X=\E_X$ gives
$\Pir_0^X C_{\IBM,\varepsilon}=\E_X[g_\varepsilon'(0)]L_t^{0}(Y)
+\tfrac12\!\int_0^t\E_X[g_\varepsilon''(|Y_s|)\mid Y]\,ds$. The first
term is zero by \cref{lem:bdy-zero}. Independence of $X,Y$ gives
$\E_X[g_\varepsilon''(|Y_s|)\mid Y]=\psi_\varepsilon(|Y_s|)$.
(2) $\Pir_{0,n}=\Pir_n^Y\Pir_0^X$ applied to (1).
(3) $\Pir_{0,0}=\E_Y\E_X=\E$. The expectation-limit step is the
statement of \cref{thm:moll-consistency} and is not re-proved here.
The interior bi-Laplacian piece is identified by
$\phi_\varepsilon^{\mathrm{int}}(r)\to\tfrac14(P_{|r|}f^{(4)})(x)$
uniformly on compacts of $(0,\infty)$ (heat-semigroup continuity)
and dominated convergence under the half-normal law of $|Y_s|$:
$\tfrac12\E_Y\!\int_0^t\phi_\varepsilon^{\mathrm{int}}(|Y_s|)\,ds\to
\tfrac18\!\int_0^t\Delta^{2}u(s,x)\,ds$. The boundary-layer piece
limit is then identified by subtraction from \cref{thm:moll-consistency}.
(4) For $n\ge 1$, $\Pir_n^Y F_\varepsilon\in\Hilb_n^Y$ is centred;
Stroock--Clark--Ocone gives the $n=1$ Malliavin-derivative
representation; the absolute-value Malliavin derivative is the
Bouleau--Hirsch standard form.
\end{proof}

\begin{openproblem}[F2: boundary-layer random functional limit]
\label{flag:bdy-layer}
\Cref{thm:pure-clock}(3) identifies the $\varepsilon\to 0$ limit of
$\E_Y\!\int_0^t b_\varepsilon(|Y_s|)\,ds$ (an expectation-level
identification of the $t^{-1/2}$ source). We make no claim about
the random-functional $L^{2}(\Omega_Y)$ limit of
$\int_0^t b_\varepsilon(|Y_s|)\,ds$ as $\varepsilon\to 0$;
\eqref{eq:b-eps-bounds} shows that the mass of $b_\varepsilon$
concentrates at scale $\varepsilon$ near $r=0$, suggesting a clock
local-time at $0$ object as the natural candidate limit, but the
identification and convergence are open. Consequently the
decomposition $\psi_\varepsilon=\phi_\varepsilon^{\mathrm{int}}+b_\varepsilon$
is used at the expectation level only (\cref{thm:pure-clock}(3));
the $(0,n)$ slices for $n\ge 1$ are recorded via the projection
formula at fixed $\varepsilon$ (\cref{thm:pure-clock}(2),(4)), not
in their $\varepsilon\to 0$ limit.
\end{openproblem}

\subsection*{The first-scenery-chaos sector $(1,\bullet)$.}

By \cref{lem:proj1-X}, the first-scenery-chaos projection of
$C_{\IBM,\varepsilon}$ satisfies
$\Pir_1^X\,C_{\IBM,\varepsilon}(f;t)=I_1^X(K_\varepsilon(\cdot\mid Y))$
with $K_\varepsilon$ given by~\eqref{eq:Keps-decomp}.

\begin{theorem}[First-scenery-chaos bidegree slices]
\label{thm:first-chaos}
For every $f\in\mathscr S(\R)$, $t\ge 0$, $\varepsilon>0$, $n\ge 0$:
\begin{enumerate}
\item (Block formula)
$\Pir_{1,n}\,C_{\IBM,\varepsilon}(f;t)=I_1^X(\Pir_n^Y K_\varepsilon(\cdot\mid Y))$.
\item (Clock-averaged block, $n=0$)
$\Pir_{1,0}\,C_{\IBM,\varepsilon}(f;t)=I_1^X(\E_Y K_\varepsilon(\cdot\mid Y))$.
This sector has zero full expectation (every $\Hilb_1^X$ element is
centred in $X$) but is generally non-zero as an element of $\Hilb_1^X$.
\item (Dual closure inheritance) The dual closure of
\cref{cor:combined} descends block-by-block: as $\varepsilon\to 0$,
$\Pir_n^Y K_\varepsilon\to\Pir_n^Y K^{\mathrm{lim}}$ in
$L^{2}(\Omega_Y;H^{-1}(\R))$, with $K^{\mathrm{lim}}=K^{\mathrm{int}}+L_t^{0}(Y)f'(x)\delta_0$.
In particular $\E_Y K_\varepsilon\to\E_Y K^{\mathrm{lim}}$ in
$H^{-1}(\R)$. The convergence is at the level of conditional
Malliavin kernels; the corresponding first-chaos objects
$\Pir_{1,n}C_{\IBM,\varepsilon}$ are ordinary
$L^{2}(\Omega)$-valued random variables at every fixed
$\varepsilon>0$, but their limit is the generalized
$Y$-conditioned first-chaos functional of
\cref{def:first-chaos-block} specified by the kernel
$\Pir_n^Y K^{\mathrm{lim}}\in L^{2}(\Omega_Y;H^{-1}(\R))$ via the
duality pairing, not an ordinary $L^{2}(\Omega)$-valued random
variable.
\end{enumerate}
\end{theorem}

\begin{proof}
(1) $\Pir_n^Y$ acts only on $\Omega_Y$ and commutes with the
$\Omega_X$-linear $I_1^X(\cdot)$:
$\Pir_n^Y[I_1^X(K_\varepsilon)]=I_1^X(\Pir_n^Y K_\varepsilon)$.
$\Pir_1^X$ is the identity on first-$X$-chaos.
(2) $n=0$ specialisation: $\Pir_0^Y=\E_Y$, full-expectation zero by
centring of $\Hilb_1^X$. Non-vanishing: the FCDC limit kernel
$K^{\mathrm{lim}}$ has explicit non-trivial $Y$-average (e.g.\ via
$-\partial_w(C(w)L_t^{w}(|Y|))$ with $\E_Y[L_t^{w}(|Y|)]$ not
identically zero).
(3) $\Pir_n^Y$ is an orthogonal projection on $L^{2}(\Omega_Y)$,
hence a contraction on $L^{2}(\Omega_Y;H^{-1}(\R))$. Apply it to
$K_\varepsilon\to K^{\mathrm{lim}}$ in $L^{2}(\Omega_Y;H^{-1}(\R))$
(\cref{cor:combined}). The $n=0$ specialisation gives the $\E_Y$
statement.
\end{proof}

\begin{remark}[Boundary local time inside $\Pir_1^X$]
\label{rem:bdy-local-time}
The boundary kernel $K_\varepsilon^{\mathrm{bdy}}\to L_t^{0}(Y)f'(x)\delta_0$
in $H^{-1/2-\delta}(\R)\subset H^{-1}(\R)$, so the first-chaos
boundary singularity is a feature of $\Pir_1^X$, distributed across
the $\Pir_{1,n}$ blocks via the $Y$-chaos expansion of $L_t^{0}(Y)$.
The expectation $\E_Y[L_t^{0}(Y)]=\sqrt{2t/\pi}$ contributes to
$\Pir_{1,0}$; the higher $Y$-chaos pieces of $L_t^{0}(Y)$
contribute to $\Pir_{1,n}$ for $n\ge 1$.
\end{remark}

\subsection*{The energy-contraction channel $(2,0)$.}

By \cref{thm:atlas} and partial expectations,
$\Pir_{2,0}\,C_{\IBM,\varepsilon}(f;t)=\E_Y[\Pir_2^X C_{\IBM,\varepsilon}(f;t)]$.

\emph{Gaussian Stroock and product formula.}
$V_r\in\Hilb_1^X$ with $V_r^{2}=\sigma_V^{2}+I_2^X(\rho_\varepsilon(r-\cdot)^{\otimes 2})$,
the second summand being the Wick-square realisation of the
second-energy contraction (we use this only as Gaussian-coordinate
bookkeeping). $W_r=I_1^X(\rho_\varepsilon'(r-\cdot))$ is first-chaos.
For $r\ge 0$ and $k\in\N$, $f^{(k)}(x+U_r)$ admits the Stroock
expansion $f^{(k)}(x+U_r)=\sum_{m\ge 0}I_m^X(h^{(k)}_{m,r,\varepsilon})$
with $h^{(k)}_{m,r,\varepsilon}=\frac{1}{m!}\E_X[f^{(k+m)}(x+U_r)]\xi_{r,\varepsilon}^{\otimes m}$.
The Gaussian product formula
$I_p^X(a)I_q^X(b)=\sum_{j=0}^{p\wedge q}j!\binom{p}{j}\binom{q}{j}I_{p+q-2j}^X(a\otimes_j b)$
(\cite[\S1.1.3]{Nualart}) isolates the chaos-2 selector $p+q-2j=2$.

\begin{proposition}[Channel enumeration]\label{prop:channels}
For every fixed $\varepsilon>0$ and $r\ge 0$,
\begin{align*}
 \Pir_2^X[f''(x+U_r)V_r^{2}]
 &=\Channel{0,0}+\Channel{2,V}+\Channel{2,\mathrm{c}}+\Channel{4},\\
 \Pir_2^X[f'(x+U_r)W_r]
 &=\Channel{1}+\Channel{3},
\end{align*}
where (with $h^{(k)}_{m,r}=h^{(k)}_{m,r,\varepsilon}$):
\begin{align*}
 \Channel{0,0}\,\,&\coloneqq\E_X[f''(x+U_r)]\cdot I_2^X(\rho_\varepsilon(r-\cdot)^{\otimes 2}),\\
 \Channel{2,V}&\coloneqq\sigma_V^{2}\cdot I_2^X(h^{(2)}_{2,r}),\\
 \Channel{2,\mathrm{c}}&\coloneqq 4\cdot I_2^X(h^{(2)}_{2,r}\otimes_1\rho_\varepsilon(r-\cdot)^{\otimes 2}),\\
 \Channel{4}&\coloneqq 12\cdot I_2^X(h^{(2)}_{4,r}\otimes_2\rho_\varepsilon(r-\cdot)^{\otimes 2}),\\
 \Channel{1}&\coloneqq I_2^X(h^{(1)}_{1,r}\,\widetilde\otimes\,\rho_\varepsilon'(r-\cdot)),\\
 \Channel{3}&\coloneqq 3\cdot I_2^X(h^{(1)}_{3,r}\otimes_1\rho_\varepsilon'(r-\cdot)).
\end{align*}
There are no further chaos-2 contributions.
\end{proposition}

\begin{proof}
Direct substitution of the Stroock expansion into $\Pir_2^X$ and
applying the product formula. The parity constraint
$p+q\equiv 2\pmod 2$ and the index bound $j\le\min(p,q)$ enumerate
all allowed $(p,q,j)$. Coefficients $j!\binom{p}{j}\binom{q}{j}$
evaluate to the listed values.
\end{proof}

\begin{theorem}[$(2,0)$ energy-contraction channel decomposition]
\label{thm:2-0}
For every $f\in\mathscr S(\R)$, $t\ge 0$, $\varepsilon>0$,
\begin{equation}\label{eq:2-0-explicit}
 \Pir_{2,0}\,C_{\IBM,\varepsilon}(f;t)
 \;=\;\sqrt{\tfrac{2t}{\pi}}\,\Pir_2^X[f'(x+U_0)V_0]
   +\tfrac12\!\int_0^t\!\E_Y\!\bigl[\Pir_2^X g_\varepsilon''(|Y_s|)\bigr]\,ds,
\end{equation}
with the second-chaos integrand given by \cref{prop:channels} as
$\Channel{0,0}+\Channel{2,V}+\Channel{2,\mathrm{c}}+\Channel{4}+\Channel{1}+\Channel{3}$
at $r=|Y_s|$. The expression lies in $\Hilb_2^X\otimes\Hilb_0^Y$ at
every $\varepsilon>0$.
\end{theorem}

\begin{proof}
By $\Pir_{2,0}=\E_Y\Pir_2^X$ applied to~\eqref{eq:moll-correction}.
For the boundary term, $g_\varepsilon'(0)$ is $Y$-independent, so
$\Pir_2^X[g_\varepsilon'(0)L_t^{0}(Y)]=L_t^{0}(Y)\Pir_2^X[f'(x+U_0)V_0]$,
and $\E_Y[L_t^{0}(Y)]=\sqrt{2t/\pi}$ by Tanaka and the half-normal
distribution of $|Y_t|$. The interior term is direct from
\cref{prop:channels} at $r=|Y_s|$.
\end{proof}

\begin{proposition}[Boundary $(2,0)$ kernel]\label{prop:2-0-bdy}
\(
 \Pir_2^X[f'(x+U_0)V_0]
 =I_2^X\!\bigl(h^{(1)}_{1,0,\varepsilon}\,\widetilde\otimes\,\rho_\varepsilon(0-\cdot)\bigr).
\)
\end{proposition}

\begin{proof}
$V_0=I_1^X(\rho_\varepsilon(0-\cdot))\in\Hilb_1^X$, so $q=1$. The
parity analysis of \cref{prop:channels} with $q=1$ allows
$p\in\{1,3\}$ with coefficients $1$ and $3$ respectively. The
$p=1,j=0$ channel contributes the displayed first term. The
$p=3,j=1$ channel is formally allowed by parity, but its 1-fold
contraction
\(
 h^{(1)}_{3,0,\varepsilon}\otimes_1\rho_\varepsilon(0-\cdot)
\)
contains the inner product
\[
 \int_\R\xi_{0,\varepsilon}(u)\,\rho_\varepsilon(0-u)\,du
 \;=\;\int_\R (D_u U_0)(D_u V_0)\,du
 \;=\;\Cov(U_0,V_0)
 \;=\;0
\]
by~\cref{lem:bdy-cov-zero} (the same identity that kills the
$\delta_0'$ boundary contribution in~\cref{thm:boundary}). Hence
the $p=3,j=1$ channel vanishes identically, leaving only the
$p=1,j=0$ term.
\end{proof}

\begin{remark}[The $(2,0)$ block is potentially $\varepsilon$-singular]
\label{rem:2-0-status}
The leading scalar channel $\Channel{0,0}$ involves the kernel
$\rho_\varepsilon(r-\cdot)^{\otimes 2}$, whose $L^{2}(\R^{2})$ norm
is $O(\varepsilon^{-1})$; multiplied by the $O(1)$ scalar
$\E_X[f''(x+U_r)]$, this channel is potentially singular as
$\varepsilon\to 0$. Whether an energy-contraction subtraction
stabilises it is part of the cross-sector Tier-A problem
(\cref{flag:summability}). We do not provide a proof of any specific
$\varepsilon\to 0$ rate for the $(2,0)$ sector; the decomposition
\cref{thm:2-0} is a fixed-$\varepsilon$ structural statement only.
\end{remark}

\cref{tab:atlas} summarises the atlas at fixed $\varepsilon$.
Three structural distinctions are worth restating:

\begin{itemize}
\item \emph{Deterministic mass vs.\ zero expectation.} Only
$\Pir_{0,0}$ has non-zero full expectation. Every other
$\Pir_{m,n}$ with $m+n\ge 1$ is centred. ``Zero expectation'' is
strictly weaker than ``zero bidegree''; the prototypical example is
$\Pir_{1,0}=I_1^X(\E_Y K_\varepsilon)$, generally non-zero
(\cref{thm:first-chaos}).
\item \emph{Two distinct boundary objects.} The $(0,0)$ boundary
residual $b_\varepsilon$ is an $X$-expectation-level object
producing the $t^{-1/2}$ source at the level of $\E$
(\cref{thm:pure-clock}(3)); the first-chaos boundary local time
$L_t^{0}(Y)f'(x)\delta_0$ is a $Y$-randomness-level object inside
$\Pir_1^X$ (\cref{thm:boundary,rem:bdy-local-time}). They are not
the same and live in different bidegree rows.
\item \emph{The $(2,0)$ contraction channel does not remove the
$(1,\bullet)$ singularity.} The six-channel decomposition is in
$\Hilb_2^X$; the obstruction is in $\Hilb_1^X$. Wick-square
renormalisation acts on $V_r^{2}-\sigma_V^{2}$ and leaves the
first-chaos rank-one $W_r$-singularity untouched
(\cref{rem:wick-not-sufficient}).
\end{itemize}

\begin{table}[h!]
\caption{Bidegree atlas of $C_{\IBM,\varepsilon}(f;t)$ at fixed
$\varepsilon$.}
\label{tab:atlas}
\footnotesize
\setlength{\tabcolsep}{4pt}
\begin{tabular}{@{}p{0.22\textwidth}p{0.32\textwidth}p{0.36\textwidth}@{}}
\textbf{Sector} & \textbf{Object} & \textbf{Status}\\[2pt]
\hline\\[-6pt]
$(0,0)$, interior of $\psi_\varepsilon$
 & bi-Laplacian shadow
 & explicit (\cref{thm:pure-clock})\\
$(0,0)$, boundary residual $b_\varepsilon$
 & singular $t^{-1/2}$ source
 & expectation limit identified; random-functional limit open (\cref{flag:bdy-layer})\\
$(0,0)$, combined
 & deterministic mass $u(t,x)-f(x)$
 & \cref{thm:moll-consistency}\\
$(0,n)$, $n\ge 1$
 & pure-clock zero-mean fluctuations
 & projection formula at fixed $\varepsilon$\\
$(1,0)$
 & clock-averaged scenery $I_1^X(\E_Y K_\varepsilon)$
 & closed in $H^{-1}$ (\cref{cor:combined})\\
$(1,n)$, $n\ge 1$
 & $Y$-resolved singular scenery channels
 & closed in $H^{-1}$ (\cref{thm:first-chaos})\\
bdy: $L_t^{0}(Y)f'(x)\delta_0$
 & first-chaos boundary local time
 & spread over $(1,n)$\\
$(2,0)$
 & six-channel energy-contraction kernel
 & explicit at fixed $\varepsilon$ (\cref{thm:2-0})\\
$(2,n)$, $n\ge 1$
 & clock-resolved second-scenery channels
 & open (\cref{flag:summability})\\
$(m,n)$, $m\ge 3$
 & higher scenery contractions
 & open (\cref{flag:summability})\\
\end{tabular}
\end{table}

Cross-sector $\varepsilon\to 0$ summability of
$\sum_{m,n}\Pir_{m,n}C_{\IBM,\varepsilon}$ in a global dual
topology would require, beyond this paper, four ingredients: a
norm on the $(2,0)$ block in which the six channels (or a suitable
contraction subtraction) converge; control of the mixed $(2,n)$
sectors; a tail estimate for $\Pir_{m,n}$ with $m\ge 3$; and
identification of the boundary residual as an $L^{2}(\Omega_Y)$
limit. None of these is needed for the fixed-$\varepsilon$ atlas
above: \cref{thm:atlas,thm:pure-clock,thm:first-chaos,thm:2-0} hold
at every $\varepsilon>0$. The four ingredients are catalogued
in~\cref{sec:open} alongside the three open problems.

%==============================================================
\section{Open problems and scope}\label{sec:open}

This section collects the three named open problems that this paper
exposes, together with directions for future work.

\subsection*{The three open problems.}

\begin{itemize}
\item[\textbf{(F1)}] \emph{Global cross-sector $\varepsilon\to 0$
summability}~(\cref{flag:summability} below). The atlas of
\cref{sec:atlas} is a fixed-$\varepsilon$ structural map. We do not
claim convergence of
$\sum_{m,n}\Pir_{m,n}C_{\IBM,\varepsilon}$ as $\varepsilon\to 0$ in
any global cross-sector dual topology. The three loci where
$\varepsilon\to 0$ is proved are: (a) the expectation
(\cref{thm:moll-consistency}); (b) the first-scenery-chaos sector
in $L^{2}(\Omega_Y;H^{-1}(\R))$ (\cref{cor:combined}); (c) the
$\Pir_{0,0}$ block, which is (a) by~\cref{thm:pure-clock}(3).

\item[\textbf{(F2)}] \emph{The boundary-layer random-functional
$L^{2}(\Omega_Y)$ limit}~(\cref{flag:bdy-layer}). The
expectation limit of $\E_Y\!\int_0^t b_\varepsilon(|Y_s|)\,ds$ is
identified as the $t^{-1/2}$ source. The random-functional limit
of $\int_0^t b_\varepsilon(|Y_s|)\,ds$ in $L^{2}(\Omega_Y)$ is open.

\item[\textbf{(F3)}] \emph{Density of an intrinsic IBM core, and
the natural home for $C_\IBM(f;t)$}~(\cref{flag:density}). The
candidate intrinsic IBM derivative is closed on its maximal pullback
core $\mathcal C_\IBM$; we are not suggesting $L^{2}(\mathcal F_\IBM)$-density
of $\mathcal C_\IBM$, and we expect it fails for the structural
reason that scenery-at-clock observables lie outside the joint
Sobolev domain in any unmollified sense. The natural home for the
unmollified pathwise correction is therefore not $L^{2}(\Omega)$
but a dual / rigged-Hilbert-space object; the first-chaos slice is
constructed in \cref{sec:dual}.
\end{itemize}

\begin{openproblem}[F1: cross-sector summability]\label{flag:summability}
For $f\in\mathscr S(\R)$ satisfying \eqref{eq:nondegen}, the
$\varepsilon\to 0$ behaviour of $\Pir_{m,n}C_{\IBM,\varepsilon}$
for $(m,n)$ with $m\ge 2$, or $m=0$ and $n\ge 1$, is open: we have
neither a uniform-in-$\varepsilon$ bound nor an identification of an
$\varepsilon\to 0$ limit in any topology. The construction of a
cross-sector dual norm in which
$\sum_{m,n}\Pir_{m,n}C_{\IBM,\varepsilon}\to C_\IBM(f;t)$ as
$\varepsilon\to 0$ remains open; it is the principal cross-sector
analytic frontier this paper exposes.
\end{openproblem}

\subsection*{Directions for future work.}

Beyond the three open problems above, the present construction
opens several directions.

\begin{enumerate}
\item \emph{Unmollified clock derivative.} Identify
$\lim_{\varepsilon\to 0}D_Y C_{\IBM,\varepsilon}(f;t)$ as a
distribution-valued kernel against $L^{2}(\R_{+})$ test integrands.
The clock component of the natural intrinsic IBM derivative is
genuinely singular at the level of local time.

\item \emph{Cylindrical reduction for $D_\IBM$.} Establish a chain
or product rule for $D_\IBM$ via a cylindrical reduction with random
representers (cf.\ \cite{Fontes-OperatorDerivative}); the natural
candidate is the joint Brownian--Malliavin representer projected
onto $\mathcal F_\IBM$.

\item \emph{Second-order intrinsic energy and PDE shadow.} Construct
$D_\IBM^{(2)}$ as a closed operator and verify that, averaged against
the law of $Y$, the second-order $\Gamma$-form built from
$D_\IBM^{(2)}$ matches the bi-Laplacian/local-time PDE shadow of
\cite{AllouZheng}. The fixed-$\varepsilon$ recovery in
\cref{thm:pure-clock}(3) gives a target on the source side; the
intrinsic-IBM-side recovery remains open.

\item \emph{Second-energy-contraction sector.} Construct the
$\Hilb_2^X$ tensor-product Bessel-potential dual block analogous
to $H^{-1}$ and study $\varepsilon\to 0$ convergence of the six
channels of \cref{thm:2-0} or a suitable contraction-subtraction
thereof.

\item \emph{Wick/Hermite hierarchy in higher chaoses.} The
$\Hilb_{m\ge 2}^X$ sectors admit contraction renormalisations
(Wick/Hermite in Gaussian coordinates). Establishing chaos-by-chaos
$L^{2}$-boundedness uniformly in $\varepsilon$, then assembling
across $(m,n)$, is a natural route towards~\cref{flag:summability}.

\item \emph{Filtration-level questions.} The natural filtration of
$\IBM$ is a separate matter; Tsirelson-type questions about whether
$\mathcal F_\IBM$ is Brownian-generated are independent of the
present construction.
\end{enumerate}

\appendix

%==============================================================
\section{Wick / Hermite renormalisation is not a stabilising renormalisation}\label{sec:wick-app}

The first-chaos obstruction (\cref{thm:var-divergence}) suggests an
obvious renormalisation: subtract the diagonal contraction
$\mathrm{Contr}_\varepsilon(r)=c_\rho^{(1)}/\varepsilon$ from the
squared first-energy direction $(\partial_r X^\varepsilon)^{2}$. In
Gaussian source coordinates this is the Wick / Hermite square in
$\Hilb_2^X$. The subtraction removes the obvious even-sector
diagonal contraction, but it neither touches the $T_3$ first-chaos
singularity (carried by $\partial_r^{2}X^\varepsilon$, a centred
Gaussian in $\Hilb_1^X$ with $\Var\sim\varepsilon^{-3}$) nor
preserves the cancellation mechanism that isolates that singularity:
the counterterm it introduces is itself a large first-chaos
$\xi_{r,\varepsilon}$-contribution. Scalar Wick-square subtraction
therefore does not remove the first-chaos obstruction and cannot by
itself provide the stabilising renormalisation of the full
correction.

The diagonal contraction is
\begin{equation}\label{eq:contraction}
  \mathrm{Contr}_\varepsilon(r)\;\coloneqq\;\E[(\partial_rX^\varepsilon_r)^{2}]
  \;=\;\frac{c_\rho^{(1)}}{\varepsilon}\quad\text{for }|r|>c\varepsilon,
\end{equation}
so Wick / Hermite centering is the Brownian-coordinate expression of
an operator-energy contraction subtraction. Define
\begin{equation}\label{eq:wick-sq}
  :(\partial_rX^\varepsilon_r)^{2}:\;\coloneqq\;
  (\partial_rX^\varepsilon_r)^{2}-c_\rho^{(1)}/\varepsilon
\end{equation}
and the renormalised correction candidate
\(
  \widehat C_{\IBM,\varepsilon}(f;t)
  \coloneqq g_\varepsilon'(0)L_t^{0}(Y)+\tfrac12\!\int_0^t\!
   :g_\varepsilon''(|Y_s|):\,ds,
\)
where $:g_\varepsilon''(r):=f''(x+U_r):(\partial_rU_r)^{2}:+f'(x+U_r)W_r$.

\begin{proposition}[Scalar Wick subtraction does not remove the first-chaos obstruction]
\label{prop:wick-insufficient}
Under the hypotheses of \cref{thm:var-divergence}, the contraction
subtraction~\eqref{eq:wick-sq} does not modify the
$\partial_r^{2}X^\varepsilon$ content of $g_\varepsilon''$, so the
first-scenery-chaos singularity in $\Pir_1^X C_{\IBM,\varepsilon}$
carried by $\partial_r^{2}X^\varepsilon$ (identified
in~\cref{thm:var-divergence}) is not removed. In addition, the
subtraction introduces a counterterm whose first-chaos contribution
is large, so the Stein cancellation that isolates the $T_3$
singularity is not preserved. Scalar Wick-square subtraction
therefore does not remove the first-chaos obstruction and cannot by
itself provide the stabilising renormalisation of the full
correction.
\end{proposition}

\begin{proof}
The subtraction~\eqref{eq:wick-sq} alters only the $V^{2}$ factor in
$g_\varepsilon''(r)=f''(x+U_r)V_r^{2}+f'(x+U_r)W_r$. The term
$f'(x+U_r)W_r$ carries the $C_\varepsilon(r)\rho_\varepsilon'(r-w)$
contribution to the first-chaos Malliavin kernel~\eqref{eq:Dkernel},
which is the source of the residual $L^{2}$-divergence in
\cref{thm:var-divergence}; it is untouched by~\eqref{eq:wick-sq}.
The subtraction further adds a counterterm
$-\tfrac{c_\rho^{(1)}}{2\varepsilon}\!\int_0^t f''(x+X^\varepsilon_{|Y_s|})\,ds$
whose first-chaos Malliavin kernel is built from the $U_r$-fibre
$\xi_{r,\varepsilon}(w)$ (via Stein on $f''(x+U_r)$). This
$\varepsilon^{-1}$-scaled $\xi_{r,\varepsilon}$-kernel modifies the
$T_1$ summand of~\eqref{eq:Dkernel}--\eqref{eq:T123} and disrupts
the cancellation~\eqref{eq:Dkernel}, while leaving
$T_3=-\partial_w(\rho_\varepsilon*C_\varepsilon L)$ intact.
\end{proof}

The renormalisation must operate at the level of a quartic composite
observable energy, not at the level of a single Gaussian square.
Two analytic routes remain open: (i) interpret
$\partial_r^{2}X$ in the unmollified limit as a Hida/Schwartz
distribution and pair $f'(X)\partial_r^{2}X$ against test
integrands in the negative-Sobolev dual; (ii) Stratonovich-regularise
the time integral $\int_0^t f'(x+X^\varepsilon_{|Y_s|})\partial_r^{2}X^\varepsilon_{|Y_s|}\,ds$
via integration by parts in the clock variable, transferring the
spatial $r$-derivative onto Brownian-clock local-time data. Neither
route is pursued here.

%==============================================================
\section{Fourier-mode PDE-shadow diagnostic}\label{sec:fourier-app}

We record a PDE-shadow diagnostic identifying, at the level of
Fourier-mode test functions, the second-order signature reproduced
by~\cref{thm:bt-correction}.

For $k\in\R$ and $t>0$, $u_k(t,x)\coloneqq\E[e^{ikZ_t}]$.
Conditioning on $|Y_t|=s$ gives $\E[e^{ikX_s}]=e^{-k^{2}s/2}$, and
$|Y_t|$ has half-normal density $p_t(s)=\sqrt{2/(\pi t)}\,e^{-s^{2}/(2t)}$
for $s\ge 0$.

\begin{proposition}[Fourier-mode bi-Laplacian shadow]
\label{prop:fourier-2nd}
Let $a=k^{2}/2$ and $q(t)\coloneqq\int_0^\infty e^{-as}p_t(s)\,ds$,
so $u_k(t,x)=e^{ikx}q(t)$. Then
\begin{equation}\label{eq:q-ode}
  q'(t)\;=\;\frac{a^{2}}{2}\,q(t)\;-\;\frac{a}{\sqrt{2\pi t}}
  \;=\;\frac{k^{4}}{8}q(t)-\frac{k^{2}}{2\sqrt{2\pi t}},
\end{equation}
and equivalently
$\partial_t u_k(t,x)=\tfrac18\Delta^{2}u_k(t,x)+\Delta f_k(x)/(2\sqrt{2\pi t})$
for $f_k(x)=e^{ikx}$.
\end{proposition}

\begin{proof}
$\phi_t(s)=(2\pi t)^{-1/2}e^{-s^{2}/(2t)}$ satisfies
$\partial_t\phi_t=\tfrac12\partial_s^{2}\phi_t$; the same holds for
$p_t=\sqrt 2\phi_t|_{s\ge 0}$ at every $s>0$. Hence
$q'(t)=\tfrac12\!\int_0^\infty e^{-as}\partial_s^{2}p_t(s)\,ds$.
IBP twice in $s$, using $\partial_s p_t(0)=0$ and $p_t(0)=\sqrt{2/(\pi t)}$:
$\int_0^\infty e^{-as}\partial_s^{2}p_t\,ds=a(-p_t(0)+aq(t))=a^{2}q(t)-a\sqrt{2/(\pi t)}$.
Dividing by 2 and substituting $a=k^{2}/2$ gives~\eqref{eq:q-ode};
the PDE follows from $\Delta e^{ikx}=-k^{2}e^{ikx}$ and
$\Delta^{2}e^{ikx}=k^{4}e^{ikx}$.
\end{proof}

Fourier-synthesis on smooth bounded $f$ gives, with $u(t,x)=\E[f(x+X_{|Y_t|})]$,
$\partial_t u(t,x)=\tfrac18\Delta^{2}u(t,x)+\Delta f(x)/(2\sqrt{2\pi t})$,
the Brownian-time PDE of~\cite{AllouZheng}.

\begin{acks}
AI tools were used as interactive assistants for drafting,
reorganizing, and generating mathematical arguments, under the
author's direction. The author is responsible for all mathematical
content and any errors.
\end{acks}

%==============================================================
\begin{thebibliography}{99}

\bibitem{AllouZheng}
H.~Allouba and W.~Zheng,
\emph{Brownian-time processes: the PDE connection and the
half-derivative generator},
Ann. Probab. \textbf{29} (2001), no.~4, 1780--1795.

\bibitem{BorodinSalminen}
A.~N.~Borodin and P.~Salminen,
\emph{Handbook of Brownian Motion --- Facts and Formulae},
2nd ed., Probability and Its Applications,
Birkh\"auser Verlag, Basel, 2002.

\bibitem{BouleauYor}
N.~Bouleau and M.~Yor,
\emph{Sur la variation quadratique des temps locaux de certaines
semimartingales},
C.~R.~Acad.~Sci.~Paris S\'er.~I Math. \textbf{292} (1981), no.~9,
491--494.

\bibitem{Fontes-OperatorDerivative}
R.~Fontes,
\emph{The Operator Derivative in Continuous Stochastic Calculus:
A Hilbert Energy Space Framework},
SSRN preprint, 2026.
Available at \url{https://papers.ssrn.com/sol3/papers.cfm?abstract_id=6381359}.

\bibitem{Janson}
S.~Janson,
\emph{Gaussian Hilbert Spaces},
Cambridge Tracts in Mathematics, vol.~129,
Cambridge University Press, Cambridge, 1997.

\bibitem{MarcusRosen}
M.~B.~Marcus and J.~Rosen,
\emph{Markov Processes, Gaussian Processes, and Local Times},
Cambridge Studies in Advanced Mathematics, vol.~100,
Cambridge University Press, 2006.

\bibitem{Nualart}
D.~Nualart,
\emph{The Malliavin Calculus and Related Topics},
2nd ed., Probability and Its Applications,
Springer-Verlag, Berlin, 2006.

\bibitem{RevuzYor}
D.~Revuz and M.~Yor,
\emph{Continuous Martingales and Brownian Motion},
3rd ed., Grundlehren der Mathematischen Wissenschaften, vol.~293,
Springer-Verlag, Berlin, 1999.

\end{thebibliography}

\end{document}
